\documentclass[12pt]{article}

\usepackage[
left=1.5in,
right=1in,
top=1in,
bottom=1in
]{geometry}
\usepackage[T1]{fontenc}
\usepackage[utf8]{inputenc}
\usepackage{textcomp}
\usepackage[strings]{underscore}
\usepackage{newunicodechar}
\usepackage{amsmath,amssymb}
\usepackage{graphicx}
\usepackage{caption}
\captionsetup{font=normalsize,labelfont=normalsize}
\usepackage{hyperref}
\usepackage{enumitem}
\usepackage{titlesec}
\titleformat{\section}[display]
  {\fontsize{14pt}{17pt}\bfseries}
  {Chapter \thesection}
  {1ex}
  {\fontsize{14pt}{17pt}\bfseries}
\titleformat{\subsection}
  {\fontsize{14pt}{17pt}\itshape}
  {\thesubsection}
  {1em}
  {}
\titleformat{\subsubsection}
  {\fontsize{12pt}{14pt}\itshape}
  {\thesubsubsection}
  {1em}
  {}


\setlength{\parindent}{0pt}
\setlength{\parskip}{0.4\baselineskip}

\begin{document}
\onehalfspacing
% This file keeps your original wording and only adds LaTeX formatting.


\begin{titlepage}
\begin{center}

{\fontsize{16pt}{19pt}\bfseries Design and Implementation of a Permissioned Blockchain for Secure Medical Record Storage \\[4pt]}


\vspace{0.8cm}

{\normalsize A Project Report submitted in partial fulfilment of the requirement for the \\[4pt]
Degree of Bachelor of Technology in Computer Science \& Engineering}

\vspace{0.8cm}

{\normalsize Maulana Abul Kalam Azad University of Technology, West Bengal}

\vspace{1.0cm}

\IfFileExists{assets/logo.png}{\includegraphics[width=3.5cm]{assets/logo.png}}{}

\vspace{0.4cm}

{\normalsize February - 2026}

\vspace{1.2cm}

\textbf{By}

\vspace{0.3cm}

Debarun Biswas

Registration No: 221000110186 of 2022-23

Examination Roll No: 10000122040

\vspace{1.2cm}

\textbf{Under the Guidance of}

\vspace{0.3cm}

Prof. Santanu Phadikar

Professor

Department of Computer Science \& Engineering

Maulana Abul Kalam Azad University of Technology

Kolkata-741249, W.B., India

\end{center}
\end{titlepage}

\newpage
\pagenumbering{roman}
\setcounter{page}{2}

\begin{center}
\textbf{Faculty of Maulana Abul Kalam Azad University of Technology}\\[0.5cm]
\IfFileExists{assets/logo.png}{\includegraphics[width=3cm]{assets/logo.png}}{}\\[0.5cm]
\textbf{CERTIFICATE}
\end{center}

\vspace{0.6cm}

This is to certify that the Dissertation Project Report entitled,
``Design and Implementation of a Permissioned Blockchain for Secure Medical Record Storage''
submitted by Debarun Biswas, University Roll No: 10000122040, Registration No: 221000110186 of 2022-23 to MAKAUT, WB, India, has been carried out under our guidance and supervision and is accepted in partial fulfilment of the requirement for the Degree of Bachelor of Technology in Computer Science and Engineering.

\vspace{3cm}
\hfill
\begin{minipage}{0.55\textwidth}

\rule{7cm}{0.4pt}

Prof. Santanu Phadikar

Professor

Department of Computer Science \& Engineering

Maulana Abul Kalam Azad University of Technology

Kolkata-741249, W.B., India

\vspace{3cm}

\rule{7cm}{0.4pt}

Mr. Mihir Sing

Head of the Department

Dept. of Computer Science \& Engineering

Maulana Abul Kalam Azad University of Technology, W.B.

\end{minipage}


\newpage
\setcounter{page}{3}

\begin{center}
\textbf{Faculty of Maulana Abul Kalam Azad University of Technology}\\[0.5cm]
\IfFileExists{assets/logo.png}{\includegraphics[width=3cm]{assets/logo.png}}{}\\[0.5cm]
\textbf{CERTIFICATE OF APPROVAL}
\end{center}

\vspace{0.6cm}

The foregoing project report is hereby approved as a creditable study of an engineering subject and presented in a manner satisfactory to warrant acceptance as prerequisite to the degree for which it has been submitted. It is understood that by this approval the undersigned do not necessarily endorse or approve any statement made, opinion expressed or conclusion drawn therein but approve the thesis only for which it is submitted.

\vspace{1.5cm}

Committee on final examination for the evaluation of the project report.

\vspace{3cm}

\begin{flushright}

\rule{7cm}{0.4pt}

Signature of the Examiner

\vspace{3cm}

\rule{7cm}{0.4pt}

Signature of the Supervisor

\end{flushright}


\newpage
\setcounter{page}{4}

\begin{center}
\textbf{DECLARATION OF ORIGINALITY AND COMPLIANCE OF ACADEMIC PROJECT REPORT}
\end{center}
\begin{center}
\includegraphics[width=3cm]{assets/logo.png}
\end{center}
\vspace{0.8cm}

I hereby declare that this project report entitled
``Design and Implementation of a Permissioned Blockchain for Secure Medical Record Storage''
contains literature survey and original research work by the undersigned candidate, as part of his Degree of Bachelor of Technology in Software Engineering.

All information has been obtained and presented in accordance with academic rules and ethical conduct.

I also declare that, as required by these rules and conduct, I have fully cited and referenced all materials and results that are not original to this work.

\vspace{1.5cm}

Name: Debarun Biswas

Examination Roll No.: 10000122040

Report Title: Design and Implementation of a Permissioned Blockchain for Secure Medical Record Storage

\vspace{2.5cm}

Date: \hspace{10cm} Signature of the candidate

\newpage
\setcounter{page}{5}

\begin{center}
\textbf{ACKNOWLEDGEMENTS}
\end{center}

\vspace{0.6cm}

First and foremost, I would like to express my sincere gratitude to my respected guides, Prof. Santanu Phadikar, Department of Computer Science \& Engineering, for providing me the opportunity to undertake this project titled “Design and Implementation of a Permissioned Blockchain for Secure Medical Record Storage.” Their constant guidance, valuable insights, and unwavering support were instrumental in the successful completion of this work. I am especially thankful for the confidence they placed in me and for their constructive feedback throughout the development of this project.

I also extend my sincere thanks to Mr. Mihir Sing for his leadership as Head of the Department during my tenure, and to all the respected faculty members of the Department of Computer Science \& Engineering for their continuous encouragement and academic support.

Finally, I am grateful to everyone who directly or indirectly contributed to the completion of this project.

\vspace{1.5cm}

Date: February 5, 2026

Place: Haringhata, West Bengal

\vspace{1cm}
\begin{flushright}
\rule{6cm}{0.4pt}

Debarun Biswas

Examination Roll No.: 10000122040

Reg. No.: 221000110186

Maulana Abul Kalam Azad University of Technology
\end{flushright}
\newpage

\section*{LIST OF FIGURES}

\noindent\hyperref[fig:3.1]{Figure 3.1 Block Structure and Hash Linking Mechanism}\dotfill\pageref{fig:3.1}\\
\hyperref[fig:3.2]{Figure 3.2 Node Interaction and Block Validation Sequence}\dotfill\pageref{fig:3.2}\\
\hyperref[fig:4.1]{Figure 4.1 Overall System Architecture}\dotfill\pageref{fig:4.1}\\
\hyperref[fig:4.2]{Figure 4.2 Protocol Interaction Diagram}\dotfill\pageref{fig:4.2}\\
\hyperref[fig:4.3]{Figure 4.3 Transaction Processing Flow Diagram}\dotfill\pageref{fig:4.3}

\newpage

\tableofcontents
\newpage
\pagenumbering{arabic}
\setcounter{page}{1}

\begin{abstract}
\fontsize{9pt}{11pt}\selectfont

This project presents the design and implementation of a custom permissioned blockchain system for securing medical record integrity. The system is built around a hybrid storage model: sensitive medical files are stored off-chain, while cryptographic hashes, metadata, validator identity, timestamps, and block linkage information are stored on a blockchain ledger. The purpose is not to replace every part of a hospital Electronic Health Record system, but to provide a tamper-evident trust layer that can prove whether a medical record has remained unchanged after it was registered.

The project has evolved beyond the earlier report. The original thesis focused mainly on a C-based blockchain core with Proof of Authority validation. The current implementation now includes three major layers: a C blockchain core, a Node.js/Express API bridge, and a Vite React/TypeScript dashboard. The C core manages blocks, transactions, hashing, Merkle root computation, Merkle proof generation, local persistence, peer-to-peer networking, synchronization, duplicate record detection, and a demonstration signature mechanism. The Express API exposes blockchain data, network health, logs, record upload endpoints, and proof-backed record verification responses to the web layer. The React frontend provides role-based pages for administrators, doctors, and patients, including dashboard statistics, block exploration, network monitoring, record viewing, record upload, activity logs, and integrity verification workflows.

The system follows a permissioned architecture suited to a healthcare consortium where participants are known entities such as hospitals, laboratories, clinics, insurance partners, or regulator-approved validator nodes. It uses a Proof of Authority style model instead of Proof of Work, because medical systems require low latency, predictable behavior, and identity-based accountability rather than energy-intensive mining. The implemented network flow also includes a three-phase-commit-inspired consensus sequence using CAN\_COMMIT, VOTE\_YES, PRE\_COMMIT, ACK\_PRE, and DO\_COMMIT messages. This gives the project a stronger distributed-systems character than the earlier version, although the current implementation remains a prototype and not a production-grade Byzantine fault tolerant system.

The unique contribution of the project lies in combining a lightweight C ledger engine with a modern hospital-facing web application. Many blockchain demonstrations are either only theoretical, only smart-contract based, or only command-line tools. This project attempts to cover the full vertical path: low-level block storage, peer networking, backend service bridge, and frontend operational dashboard. It also uses off-chain storage and on-chain hashes, which is important for healthcare privacy because raw patient data should not be replicated across every node in a blockchain network.

At the same time, the current system has limitations. The signature layer is presently a mock demonstration signature rather than real RSA/ECDSA/Ed25519 verification. The backend authentication middleware exclusively supports mock user IDs (via the X-Mock-User-ID header) to provide a streamlined, zero-dependency demonstration mode, completely removing any reliance on external services like Supabase. The upload flow is now fully integrated end to end: the API receives the file, calls the C core ENCRYPT command to securely encrypt the record under doctor and patient envelopes in the node's off-chain storage, and then calls the C viewer ADD command to commit the transaction to the ledger, returning the block index, SHA-256 hash, and data pointer to the frontend. The verify flow is also integrated: the API passes the uploaded file to the C viewer VERIFY\_FILE command, which recomputes the uploaded file hash, uses the transaction hash index to locate the matching block and transaction, and returns MATCH or NOT\_FOUND. For a match, it returns transaction metadata, block header fields, and a Merkle proof; the frontend independently reconstructs the leaf and proof root before accepting the result. Transport security, production key management, real header-chain validation, database hardening, rate limiting, test coverage, and real healthcare compliance controls remain future work. These limitations do not invalidate the academic value of the project; rather, they define the boundary between the current prototype and a production-ready healthcare infrastructure system.


\vspace{1cm}\noindent\textbf{Keywords:} Permissioned Blockchain, Proof of Authority, Medical Records, SHA-256, RSA, Distributed Ledger, Decentralized Healthcare.
\end{abstract}

\newpage
\section{INTRODUCTION}

\subsection{Background}

Healthcare data is one of the most sensitive classes of digital information. A medical record may contain diagnosis history, prescriptions, laboratory reports, imaging results, allergies, surgical notes, insurance references, and long-term patient identifiers. The correctness of this data directly affects patient safety. If a medical record is altered, deleted, forged, or silently replaced, the consequence may not be only financial loss or administrative confusion; it can affect clinical decisions.

Most healthcare institutions still depend on centralized Electronic Health Record systems. These systems are efficient for storage and retrieval, but they place a large amount of trust in one database, one administrative domain, or one privileged operator group. If that trusted center is compromised, the stored record and the audit trail may both become unreliable. A database administrator, malware operator, insider attacker, or stolen privileged credential can potentially change records and logs in the same system.

Blockchain technology provides a different approach to integrity. A blockchain is an append-only sequence of blocks where each block contains a cryptographic hash of the previous block. If an older block is modified, its hash changes. This breaks the link to the next block and allows tampering to be detected. For medical data, this does not mean raw records should be placed directly on-chain. Raw medical data must remain confidential and must be handled under strict legal and ethical constraints. Therefore, this project uses a hybrid architecture: store the sensitive medical file off-chain, store only the hash and metadata on-chain, and use the blockchain as an integrity anchor.

The project is a custom implementation rather than a direct use of Ethereum, Hyperledger Fabric, or another blockchain platform. This makes it academically valuable because the internal mechanisms are visible: block creation, hash computation, block verification, disk persistence, socket networking, synchronization, and frontend presentation are all implemented in the repository. The project is not just a smart contract or a database wrapper. It is a full stack prototype that shows how a medical record integrity system can be built from the systems layer upward.

\subsection{Problem Statement}

The central problem addressed by this project is:

How can a healthcare system maintain a tamper-evident, auditable, and distributed record of medical file integrity without exposing sensitive medical data directly on a blockchain?

This problem can be divided into several sub-problems:

1. How can a medical file be represented on-chain without storing the file itself?
2. How can block history be protected against undetected modification?
3. How can validator identity be represented in a permissioned network?
4. How can duplicate records or replayed file hashes be detected?
5. How can multiple local nodes synchronize blockchain state?
6. How can the blockchain core be exposed to a web dashboard?
7. How can administrators, doctors, and patients view or interact with the system?
8. What limitations must be resolved before production use?

The project currently answers these questions at prototype level. It implements the core ledger and dashboard concepts, while leaving several security-hardening and production-integration tasks for future work.

\subsection{Objectives and Scope}

The main objectives of the project are:

The primary objective of this project is to build a custom blockchain core for medical record integrity. This involves storing medical file hashes on-chain while keeping the sensitive files off-chain, and providing block hash chaining using SHA-256. The system aims to provide validator identity fields and demonstration block signatures, support local disk persistence, and enable peer-to-peer node communication over TCP with synchronization capabilities. Additional objectives include implementing duplicate transaction detection, building an Express API bridge, and developing a React dashboard for users to inspect the chain and network health. Finally, the project documents its strengths, weaknesses, and future improvements.

The scope is intentionally focused on integrity rather than full hospital data management. The system does not attempt to replace an EHR. It does not implement all clinical workflows. It does not store raw medical content on-chain. It is best understood as a tamper-evident integrity and audit layer that could sit beside existing hospital systems.


\newpage
\section{EXHAUSTIVE ANALYSIS AND COMPARATIVE EVALUATION}

Exhaustive Analysis of a Full-Stack Permissioned Blockchain Architecture for Medical Record Integrity: A Comparative Evaluation Against Contemporary Decentralized Healthcare Paradigms

\subsubsection{Introduction to the Centralized Healthcare Data Crisis and the Distributed Ledger Imperative}
The global digitization of healthcare and the universal adoption of Electronic Health Records (EHRs) have catalyzed unprecedented advancements in care coordination, longitudinal patient tracking, and clinical decision-making. However, the architectural foundations of these systems rely almost exclusively on centralized relational databases and siloed cloud repositories. This structural centralization introduces severe, systemic vulnerabilities. Traditional EHR systems suffer from inherent single points of failure, rendering them highly susceptible to catastrophic network outages, sophisticated ransomware attacks, and imperceptible data manipulation by privileged actors ~\cite{1}. In a conventional hospital database architecture, the entity responsible for storing sensitive medical files is identical to the entity managing the system's access logs and audit trails. Consequently, a compromised credential or a malicious database administrator possesses the capability to alter a patient's diagnostic history while simultaneously erasing the forensic evidence of that alteration, thereby directly jeopardizing patient safety and institutional liability ~\cite{4}.

To mitigate the fragility of centralized trust brokers, the medical informatics community has aggressively explored distributed ledger technologies (DLTs) and blockchain architectures. By leveraging cryptographic hashing, decentralized consensus mechanisms, and append-only immutable data structures, blockchains theoretically provide a framework for establishing trust, provenance, and tamper-evidence without relying on a centralized administrative domain ~\cite{6}. Despite this theoretical promise, the direct application of generalized, public blockchains—such as the Ethereum mainnet or Bitcoin—to healthcare environments has introduced debilitating operational limitations. These limitations manifest as extortionate transaction gas costs, unmanageable network latency, severe computational bloat, and profound conflicts with stringent global privacy regulations such as the Health Insurance Portability and Accountability Act (HIPAA) and the General Data Protection Regulation (GDPR) ~\cite{6}.

This comprehensive report provides a deep, expert-level evaluation of a newly proposed, custom-built permissioned blockchain system designed explicitly for medical record integrity. To establish a rigorous context for this evaluation, the analysis systematically dissects a massive corpus of over twenty seminal and contemporary academic frameworks within the blockchain healthcare domain. By interrogating the technical, cryptographic, and operational paradigms of legacy systems, this report illuminates the persistent limitations that have plagued the industry for the past decade. Subsequently, the narrative meticulously details how the proposed full-stack architecture—comprising a custom C-based ledger core, a Node.js Express API middleware bridge, and a Vite React TypeScript frontend dashboard—resolves these historical limitations, thereby charting a highly pragmatic and secure trajectory for the future of decentralized medical informatics.

\subsubsection{Architectural Deconstruction of the Proposed Medical Record Integrity System}
Before engaging in a comparative analysis with existing academic literature, it is imperative to establish a precise technical understanding of the proposed project's architecture. The system under evaluation departs significantly from the heavy abstractions of commercial enterprise blockchain frameworks and public smart contract networks. Instead, it is engineered as a highly integrated, full-stack prototype structured across three distinct operational layers ~\cite{4}.

The foundation of the architecture is a low-level C blockchain core. By bypassing virtual machines entirely, the core ledger engine executes natively, managing fixed-size C structures for blocks and transactions. This layer is responsible for the deterministic computation of SHA-256 hashes, the generation and verification of Merkle roots, peer-to-peer (P2P) networking over raw TCP sockets, and the persistent binary storage of the append-only ledger via a blockchain.dat file ~\cite{4}. To guarantee data durability, the system enforces explicit disk synchronization operations, utilizing fsync on Unix-like environments and \_commit on Windows, ensuring that newly appended blocks are physically written to the storage medium rather than languishing in volatile memory caches ~\cite{4}. Furthermore, the C core executes complex state synchronization mechanisms, allowing lagging nodes to request missing blocks via GET\_BLOCK sequences after comparing network chain heights ~\cite{4}.

Sitting above the foundational C core is the Node.js and Express API bridge. Because modern web browsers cannot securely execute native C binaries or directly read binary ledger files from a local file system, this middleware layer acts as a vital translation matrix. The Express API utilizes child\_process.exec to invoke the C core's command-line viewer tools, parses the resulting raw text outputs, and serializes the data into structured JSON payloads suitable for web consumption ~\cite{4}. Additionally, the API layer actively scans designated local ports to monitor the health of the P2P network, aggregates system logs for auditability, and manages a localized, zero-dependency mock authentication system via the X-Mock-User-ID header, providing a streamlined demonstration mode free from external cloud dependencies ~\cite{4}.

The uppermost layer of the architecture is the React and TypeScript frontend dashboard, powered by the Vite build tool. This layer translates complex cryptographic operations into an intuitive, role-based graphical user interface tailored for hospital administrators, clinical doctors, and patients ~\cite{4}. The dashboard features comprehensive block exploration interfaces, network latency visualizations, aggregated activity logs, and a deeply integrated medical record upload and verification workflow ~\cite{4}. By surfacing the underlying ledger mechanics through a modern, hospital-facing web application, the proposed system transcends the theoretical limitations of command-line-only academic prototypes, providing a tangible operational face to the decentralized network.

Crucially, the system employs a strict hybrid storage paradigm. It recognizes that raw protected health information (PHI) must never be permanently replicated across a distributed ledger. When a medical document is uploaded, the Node.js API invokes the C core's ENCRYPT command, which currently executes a mock envelope encryption sequence (slated for an upgrade to AES-256-GCM in production) to secure the payload ~\cite{4}. The resulting encrypted .enc file is stored securely off-chain in a localized offchain/records directory on the specific node ~\cite{4}. The blockchain itself merely stores a lightweight metadata payload containing the patient identifier, the doctor identifier, a pointer to the off-chain file, a timestamp, and, most importantly, the SHA-256 cryptographic hash of the encrypted document ~\cite{4}.

\subsubsection{Resolving the Centralized Database Crisis, Insider Threats, and AI-Driven Cyberattacks}
The academic literature universally identifies the centralized nature of contemporary EHR systems as their most critical failing. In a comprehensive review of EHR vulnerabilities, researchers note that centralized structures consolidate the entire patient database onto a single server, creating an immense target for external attackers and a catastrophic single point of failure ~\cite{1}. If the central server experiences hardware failure or a targeted distributed denial-of-service (DDoS) attack, the entire hospital network loses access to life-saving clinical histories ~\cite{1}.

Beyond external attacks, the threat of malicious insiders—authorized personnel who manipulate or exfiltrate data—remains a pervasive challenge. Hurst et al. (2022) emphasize that the digitization of patient records inherently exposes them to privilege abuse, proposing a supervised machine learning anomaly detection approach utilizing Support Vector Machines (SVM) to flag unauthorized access behaviors ~\cite{12}. Similarly, Al-Shehari (2021) proposed a multi-tiered framework for insider threat prevention, noting that traditional Role-Based Access Control (RBAC) models often fail to prevent privilege escalation or unauthorized data duplication ~\cite{1}.

Attempting to secure collaborative, multi-institutional environments, Paracha et al. (2024) introduced the Advanced Privacy-Preserving Decentralized Federated Learning (APPDFL) framework ~\cite{13}. This system sought to detect insider threats across healthcare institutions using a peer-to-peer federated learning architecture enhanced with differential privacy, theoretically eliminating the centralized server bottleneck ~\cite{13}. However, federated learning models themselves face profound vulnerabilities. As Rani et al. (2026) highlight in their research on AI-driven threat propagation and zero-day cloud vulnerabilities, machine learning models can be subject to data poisoning if the provenance of the underlying training data is not cryptographically anchored ~\cite{15}. Samant et al. (2025) further explored this in IoT-enabled healthcare environments, introducing the "SecHealth" framework to mitigate internal and external malicious behaviors at the edge ~\cite{17}. Despite these highly sophisticated AI and edge-computing advancements, the fundamental reliance on centralized audit logs leaves the historical integrity of the medical record in question; an attacker who gains root access can simply rewrite the AI's training logs or delete the anomaly detection alerts.

The proposed system resolves this foundational crisis not through probabilistic AI anomaly detection, but through deterministic cryptographic immutability. By decentralizing the trust broker, the project entirely eradicates the single point of failure ~\cite{3}. In the proposed architecture, when a doctor uploads a medical record, the transaction must be broadcast to a peer-to-peer network of validator nodes ~\cite{4}. The system employs a customized, three-phase-commit-inspired (3PC) network flow consisting of CAN\_COMMIT, VOTE\_YES, PRE\_COMMIT, ACK\_PRE, and DO\_COMMIT messages ~\cite{4}. This synchronous agreement mechanism ensures that a single malicious database administrator cannot unilaterally alter a patient's history.

If an insider threat actor circumvents hospital access controls and directly modifies a raw medical file residing in the offchain/records directory, the cryptographic hash of that altered file will immediately change. Because the original, untampered SHA-256 hash is permanently locked inside the blockchain's Merkle root—which is subsequently hashed into the block header and linked to all subsequent blocks via the previous\_hash parameter—the tampering becomes mathematically evident ~\cite{4}. When a subsequent audit or clinical review triggers the VERIFY\_FILE command, the system will recompute the hash of the tampered off-chain file, compare it to the immutable on-chain index, and immediately return a failure ~\cite{4}. This paradigm shift transforms the healthcare network from a reactive, perimeter-defense model into a proactive, mathematically verifiable tamper-evident ecosystem.

\subsubsection{Circumventing the Public Blockchain Bottleneck: Cost, Latency, and EVM Constraints}
The initial wave of blockchain healthcare research focused heavily on leveraging existing public networks, predominantly Ethereum, to secure medical records. The most seminal of these works was MedRec, introduced by Azaria et al. (2016) ~\cite{19} (A. Ekblaw et al., "MedRec: Using Blockchain for Medical Data Access..."). MedRec proposed a decentralized record management system utilizing Ethereum smart contracts to manage authentication, confidentiality, and data sharing across disparate treatment sites ~\cite{20}. To sustain the network, MedRec innovatively proposed incentivizing medical stakeholders—such as government-funded researchers and public health authorities—to act as network "miners" executing Proof of Work (PoW) algorithms, rewarding them with access to anonymized, aggregate epidemiological data ~\cite{20}.

While MedRec was conceptually groundbreaking, its reliance on a public ledger and PoW introduced severe operational and economic bottlenecks. As the Ethereum network matured, the transaction latency and variable execution costs (gas fees) rendered it entirely unviable for high-throughput hospital environments ~\cite{21}. Agoro et al. (2025) and Huang et al. (2025) corroborate this, demonstrating that existing public blockchains struggle with transaction costs to such a degree that updating patient consent models cannot keep pace with the real-time demands of clinical users ~\cite{10}.

The economic limitations of public Ethereum for healthcare are rooted deeply in the architecture of the Ethereum Virtual Machine (EVM). A comprehensive systematic review of blockchain in EHRs details how developers face immense cost-optimization limitations due to the EVM allocating memory in rigid 256-bit words ~\cite{6}. To avoid bankrupting a hospital system with gas fees, developers must employ highly complex, unnatural programming paradigms. These include "packing variables" sequentially to optimize memory, utilizing "Boolean boxing" to aggregate multiple boolean values into a single word via bitwise operations, avoiding dynamic arrays in favor of mappings, and strictly limiting string lengths to under 32 bytes ~\cite{6}. Attempting to manage the vast, dynamic datasets of a hospital network within these suffocating computational constraints inevitably leads to severe network congestion and systemic failure ~\cite{6}. Furthermore, while Layer-2 rollups have been proposed to reduce transaction costs, they introduce secondary complexities regarding data availability and settlement finality ~\cite{9}.

The proposed architecture entirely circumvents the EVM cost constraints and the public blockchain bottleneck by executing a custom-built, native C blockchain core ~\cite{4}. Because the system is engineered as a private, permissioned ledger tailored for a specific healthcare consortium, there is no cryptocurrency, no competitive gas market, and absolutely no need to employ convoluted Solidity optimizations to save computational bytes ~\cite{4}.

The architecture enforces deterministic memory allocation by defining rigid, fixed-width C struct representations for blocks and transactions ~\cite{4}. This design eliminates the garbage collection overhead and dynamic memory reallocation latencies that plague higher-level virtual machines. The C core writes these binary structures directly to the disk, ensuring ultra-fast append operations that can easily accommodate the high-throughput requirements of a busy hospital network. By dropping the Ethereum dependency entirely, the proposed project achieves predictable, zero-cost computational performance, allowing hospital IT budgets to scale based on standard hardware provisioning rather than volatile cryptocurrency gas markets.

\subsubsection{Deconstructing Smart Contracts, Access Control, and Semantic Interoperability}
To manage the complex permissions required by HIPAA and GDPR, subsequent researchers developed extensive smart contract logic to govern access control on the blockchain. Ancile, proposed by Dagher et al. (2018), utilized Ethereum-based smart contracts combined with advanced cryptographic obfuscation techniques to provide secure interoperability and highly granular access control for patients, healthcare providers, and third-party researchers ~\cite{23} (Dagher et al., "Ancile: Privacy-Preserving Framework..."). While Ancile successfully modeled complex permissions theoretically, its reliance on heavy smart contract execution inherited the computational bloat discussed previously ~\cite{24}.

Addressing the semantic interoperability problem, FHIRChain (Zhang et al., 2018) attempted to encapsulate the standard HL7 Fast Healthcare Interoperability Resources (FHIR) data models directly within a blockchain architecture ~\cite{26}. FHIRChain utilized digital health identities—functioning as public/private cryptographic key pairs—to authenticate clinical participants for collaborative, remote cancer care decision-making ~\cite{26}. However, forcing complex, deeply nested FHIR JSON objects into on-chain smart contract state variables introduces significant parsing overhead and throughput bottlenecks ~\cite{28}. More recent approaches, such as integrating domain-specific Large Language Models (LLMs) like ClinicalBERT to automatically extract semantic attributes for Attribute-Based Encryption (ABE), push the computational boundaries of smart contracts even further ~\cite{30}.

The proposed system deliberately rejects the trend of placing complex business logic and semantic data models onto the blockchain layer. Recognizing that a ledger should be a mathematically pristine engine of integrity rather than a bloated application server, the project keeps the on-chain data footprint aggressively lightweight ~\cite{4}. The transaction struct only stores foundational metadata: patient\_id, doctor\_id, data\_hash, data\_pointer, and a timestamp ~\cite{4}.

The heavy lifting of semantic interoperability, role-based access control (RBAC), and user session management is entirely decoupled from the C core and handled by the Node.js Express API and the React frontend ~\cite{4}. This separation of concerns means the blockchain never has to parse an HL7 FHIR object or execute a ClinicalBERT inference model; it merely verifies the cryptographic integrity of the pointers. This architectural decision achieves the horizontal scalability that frameworks like FHIRChain struggle to attain, allowing the web layer to iterate rapidly on UI/UX and complex business logic without requiring dangerous hard forks of the underlying consensus rules ~\cite{4}.

\subsubsection{The Hybrid Storage Dilemma and the Fallacy of IPFS Persistence}
It is a universally accepted maxim in decentralized healthcare informatics that raw medical data must never reside directly on a blockchain. Doing so violates the right to erasure (data destruction), drastically inflates the size of the ledger, and exposes highly sensitive information to network-wide replication ~\cite{4}. Consequently, hybrid storage paradigms emerged as the industry standard.

MeDShare (Xia et al., 2017) proposed a trust-less data-sharing scheme among cloud service providers utilizing a consortium blockchain ~\cite{32}. In this model, the raw data remained entrenched within traditional, third-party cloud repositories (e.g., AWS, Azure), while the blockchain was relegated to managing data provenance, access auditing, and user control ~\cite{34}. While effective at keeping the blockchain lightweight, MeDShare maintained a heavy dependency on external corporate cloud infrastructure, undermining the ethos of true decentralization.

OmniPHR (Roehrs et al., 2017) proposed a distributed architecture model to integrate Personal Health Records (PHRs) via a routing overlay network, physically dividing health records into discrete data blocks managed by high-capacity superpeers ~\cite{35}. Dubovitskaya et al. (2017/2018) presented a framework tailored specifically for radiation oncology, utilizing a permissioned blockchain alongside both local and cloud databases, relying on symmetric keys and Uniquely Identifiable Information (UIIP) hashes to retrieve the off-chain data securely ~\cite{37}.

In recent years, the InterPlanetary File System (IPFS) has become the default off-chain storage mechanism for academic prototypes. EHRChain (2022) utilized a dual-blockchain approach based on Hyperledger Sawtooth combined with IPFS for distributed file storage, storing AES-encrypted records on IPFS and anchoring the resultant hash references on-chain ~\cite{39}. Similarly, the BiiMED framework (Bouchkaren / Kung et al., 2021) and architectures analyzed by Pineda Rincón et al. (2021) heavily rely on private Ethereum smart contracts interacting with IPFS and a Trusted Third Party Auditor (TTPA) to guarantee data integrity ~\cite{41}.

However, the reliance on IPFS introduces severe, often overlooked operational risks. IPFS is a decentralized swarm; it does not guarantee permanent data storage. If a file is not actively "pinned" by a node on the network, it can be lost to garbage collection ~\cite{44}. Furthermore, retrieving large medical imaging files (e.g., DICOM scans) across an IPFS swarm introduces unpredictable latency, making it highly suboptimal for emergency clinical environments where immediate access to a patient's history is a matter of life and death ~\cite{46}.

The proposed system brilliantly bypasses the latency, data-loss risks, and third-party dependencies of IPFS and cloud providers by implementing a natively managed, deterministic local off-chain storage architecture ~\cite{4}. When the React frontend initiates a file upload, the Node.js API receives the payload and triggers the C core's ENCRYPT command ~\cite{4}. The C core executes mock envelope encryption (securing the payload utilizing keys derived from the patient and doctor identifiers) and immediately commits the resulting .enc file to a dedicated offchain/records directory physically located on the node's local disk ~\cite{4}.

This guarantees absolute high availability and deterministic retrieval speeds, as the file resides on controlled, HIPAA-compliant hospital infrastructure rather than an ephemeral public swarm. By keeping the cryptographic encryption and storage routines tightly coupled within the C core, the proposed system achieves the privacy benefits of off-chain storage without sacrificing the reliability required by critical healthcare infrastructure.

\subsubsection{Evolutionary Consensus Models: Transitioning to Proof of Authority}
The environmental and computational waste generated by Proof of Work (PoW) consensus algorithms rendered early systems like MedRec functionally obsolete for enterprise use. To bypass these inefficiencies, researchers began exploring alternative consensus mechanisms. MedBlock (Fan et al., 2018) utilized a Delegated Proof of Stake (DPoS) algorithm to select proxy nodes responsible for re-encrypting EMRs, thereby achieving consensus without massive energy consumption or network congestion ~\cite{48}. However, stake-based systems (PoS/DPoS) introduce complex tokenomics and the risk of network monopolization by wealthy nodes, an inappropriate dynamic for a collaborative hospital consortium.

Consequently, the academic consensus has coalesced around Proof of Authority (PoA). Various recent studies by Yaqoob, Arul, and Shaik (2020-2025) confirm that PoA algorithms—which rely on the verified identity and legal reputation of pre-approved validators rather than computational hashing power or financial stake—dramatically reduce latency and computational overhead in healthcare networks ~\cite{50}. The BlockCare project (2024) successfully utilized PoA to eliminate the risks of central databases while ensuring the tamper-proof storage of digital health records ~\cite{54}. Research into healthcare asset tracking further highlights that PoA provides superior tracking accuracy, highly efficient gas usage (when mapped to an EVM), and robust data integrity suitable for sensitive environments ~\cite{53}.

The proposed system adopts this paradigm shift but implements it without the heavy infrastructural requirements of platforms like Hyperledger Fabric ~\cite{4}. The system employs a custom-coded, three-phase-commit-inspired (3PC) consensus protocol over raw TCP sockets ~\cite{4}.

\subsubsection{When a node proposes a new block containing a medical record transaction, the following deterministic sequence occurs:}
The proposer broadcasts a CAN\_COMMIT message containing the serialized block to all known peer ports.
Peers receive the block, independently verify the Merkle root, the previous hash linkage, and the signature validity, and reply with a VOTE\_YES message if the block is structurally sound.
Upon receiving a majority of unique votes (preventing duplicate vote spoofing via a vote mask), the proposer broadcasts a PRE\_COMMIT message.
Peers acknowledge their readiness via ACK\_PRE.
Finally, the proposer appends the block to its local blockchain.dat file and broadcasts a DO\_COMMIT message, instructing all peers to finalize the append operation locally ~\cite{4}.

This PoA implementation entirely eliminates mining waste and tokenomic complexities. Because the validators in a hospital consortium are known, legally bound entities, this synchronous voting mechanism ensures rapid, deterministic finality, perfectly aligning with the high-throughput demands of medical informatics.

\subsubsection{Cryptographic Identity Resolution and O(1) Duplicate Record Prevention}
A persistent and highly dangerous operational challenge within modern EHR systems is the proliferation of duplicate medical records. When a patient's history is fragmented across multiple duplicate profiles, clinicians lack a comprehensive view of allergies, medications, and prior diagnoses, leading to severe clinical errors. Tayal \& Murumkar (2026) highlight that poor patient identity protection and duplicate records directly impact data accuracy, necessitating repeated tests and undermining confidence in digital health solutions ~\cite{56} (Tayal \& Murumkar, "Patient Identity Protection and Duplicate Record Prevention..."). They emphasize the need for robust Master Patient Indexes (MPI) and biometrics.

Adapa (2025/2026) explored Master Data Management (MDM) frameworks that leverage blockchain's immutability to mitigate redundancy and improve traceability across enterprise data ecosystems ~\cite{58}. Khurshid et al. (2020) also tested specific blockchain applications for patient identity management ~\cite{60}. However, the vast majority of these MDM frameworks rely on probabilistic semantic matching algorithms—comparing names, birth dates, and addresses—which are computationally expensive and prone to false positives.

The proposed system approaches the duplicate record problem not through heuristic semantic matching, but through mathematically precise cryptography. It implements a transaction hash index engineered with a pure-C Bloom filter utilizing three independent DJB2-based hash functions ~\cite{4}. This Bloom filter acts as a lightning-fast, constant-time O(1) pre-filtering cache at the ingestion layer.

When a doctor attempts to upload a medical document, the system computes the SHA-256 hash of the file. Before engaging the complex consensus network, the system checks the Bloom filter. If the hash triggers a positive match, the system queries an in-memory separate chaining hash map to confirm the duplication against the existing ledger ~\cite{4}. If an exact duplicate is found, the system instantly rejects the upload with an HTTP 409 Conflict response via the Express API, stopping the redundancy before it ever hits the blockchain ~\cite{4}.

Furthermore, to ensure this capability scales to enterprise levels, the project report details an architectural design for LevelDB Log-Structured Merge (LSM) tree integration. This proposed LevelDB schema will allow the system to scale duplicate detection to over 1 billion medical records while maintaining O(log N) SSD lookup speeds and capping physical RAM usage at approximately 2 GB, a feat impossible with pure in-memory indexes ~\cite{4}. This provides a definitive, cryptographically precise solution to file duplication that semantic MDM systems struggle to match.

\subsubsection{Redefining the Light-Client Trust Boundary Through Zero-Trust Browser Verification}
A critical, yet rarely discussed, vulnerability in many proposed web-based blockchain dashboards is the light-client trust boundary. Many academic prototypes provide a web interface that simply queries a backend API; if the API says a record is valid, the frontend displays a green checkmark. This architecture is fundamentally flawed. It merely shifts the single point of failure from the database to the web server middleware. If the backend API is compromised by an attacker, it can simply lie to the frontend, claiming that a tampered medical file is cryptographically valid.

The proposed project resolves this trust boundary vulnerability by implementing advanced, zero-trust Merkle proof verification directly within the browser ~\cite{4}. This represents a significant advancement over typical web-blockchain bridges.

\subsubsection{The verification workflow operates as follows:}
A patient or auditor uploads a medical file to the React frontend for verification.
The Node.js API passes the file to the C core's VERIFY\_FILE command.
The C core recomputes the SHA-256 hash, locates the transaction via the index, and generates a full Merkle proof—containing the transaction leaf hash, all necessary sibling hashes, and the left/right positional directions required to climb the tree ~\cite{4}.
The API returns this Merkle proof, along with the block header metadata, to the frontend.
Crucially, the React frontend does not accept a MATCH response from the API at face value. The browser independently recomputes the SHA-256 hash of the uploaded file locally using Web Crypto APIs.
The frontend reconstructs the canonical transaction leaf from the returned metadata, mathematically applies the Merkle proof sibling hashes step-by-step up the tree, and verifies that the computed root exactly matches the Merkle root provided in the block header ~\cite{4}.
Finally, it recomputes the block hash from the header fields to ensure total structural integrity.

This zero-trust architecture ensures that even if the Express API middleware is completely compromised by an insider threat, it cannot mathematically spoof a valid cryptographic proof for an altered medical file. The browser acts as an independent cryptographic auditor, providing true end-to-end integrity verification for the end-user.

\subsubsection{Comprehensive Synthesis and Future Production Trajectory}
The proposed Permissioned Blockchain-Based Medical Record Integrity System represents a formidable, highly pragmatic leap forward in the application of distributed ledger technology to healthcare informatics. By meticulously synthesizing the lessons learned from a decade of academic literature—ranging from the foundational public-chain concepts of MedRec and FHIRChain to the hybrid storage models of MeDShare and OmniPHR—the system intentionally strips away the unnecessary computational bloat of generalized smart contract platforms.

By executing a highly deterministic, low-level C blockchain core combined with a 3PC Proof of Authority consensus protocol, the architecture entirely eradicates the extortionate gas fees, memory constraints, and latency bottlenecks associated with the Ethereum Virtual Machine. Furthermore, its strict adherence to a localized, off-chain encrypted storage paradigm ensures maximum patient confidentiality and HIPAA compliance without relying on the volatile data availability risks inherent in IPFS swarms or third-party cloud monopolies.

Most importantly, the system transcends the theoretical boundaries of academic whitepapers by delivering a fully integrated, vertically demonstrable full-stack application. Through advanced cryptographic implementations—such as O(1) Bloom filter duplicate pre-filtering, rigid C-struct memory management, and zero-trust, browser-side Merkle proof validation—the framework definitively solves the entrenched problems of centralized data corruption, insider threats, and fragmented patient identities.

While the author correctly identifies that the system is currently a high-fidelity prototype, the roadmap to production is clear and well-defined ~\cite{4}. The transition from mock XOR encryption to military-grade AES-256-GCM, the upgrading of string-based signatures to true Elliptic Curve (ECDSA) or post-quantum cryptography (PQC) 15, and the encapsulation of plain TCP node traffic within mutual TLS (mTLS) 4 will elevate this architecture to enterprise-grade security. The architectural foundation evaluated in this report successfully maps out the most secure, highly performant, and economically viable trajectory for the future of decentralized electronic health record integrity.


\subsubsection{Summary Comparison}

\begin{table}[h!]
\centering
\resizebox{\textwidth}{!}{%
\begin{tabular}{|l|l|l|l|l|}
\hline
System & Verification & Storage & Limitation & Proposed Improvement \\ \hline
MedRec [19] (A. Ekblaw et al., "MedRec: Using Blockchain for Medical Data Access...") & Full chain & Off-chain & Slow & Merkle proof \\ \hline
Ancile [22] (Dagher et al., "Ancile: Privacy-Preserving Framework...") & Smart contract & On-chain heavy & High latency & Lightweight core \\ \hline
IPFS systems [38] (EHRChain, "Addressing the Challenges of Electronic Health Records Using Blockchain and IPFS") & Hash-based & IPFS & Availability risk & Local deterministic storage \\ \hline
Traditional MDM [56] (Tayal \& Murumkar, "Patient Identity Protection and Duplicate Record Prevention...") & Probabilistic & Relational DB & Computational cost & O(1) Bloom Filter \\ \hline
**Proposed System** & **Zero-trust Browser** & **Local Deterministic** & **-** & **Full-Stack Integrated** \\ \hline


\newpage
\section{PROPOSED SYSTEM}

\subsection{System Overview}

\end{tabular}
}
\end{table}
The project is a full stack prototype of a permissioned medical blockchain system. It has three primary components.

\subsection{C Blockchain Core}

The C core is located mainly under:

backend/core/src

It defines the block and transaction structures, computes SHA-256 hashes, stores blocks in a local binary ledger file, provides a viewer tool, validates the chain, and runs networked node instances. It also stores key files and off-chain encrypted record files under the core directory or demo instance folders.

\subsection{Express API}

The API is located under:

backend/api

It uses Node.js and Express. Its role is to expose blockchain information to the frontend. It runs the C viewer binary through child processes, parses the text output, and returns JSON. It also scans node ports to report network health and reads log files from demo node instances.

\subsection{React Frontend}

The frontend is located under:

frontend

It is a Vite React TypeScript application. It provides pages for:

\begin{itemize}[leftmargin=1.8em]
\item Login
\item Dashboard
\item Block Explorer
\item Network Health
\item System Logs
\item Upload Record
\item My Records / All Records
\item Verify Record
\item Patient Search
\item Activity Logs
\item Settings
\end{itemize}

The frontend gives the project an operational face. Instead of only being a command-line blockchain, the project can be demonstrated as a hospital-style web dashboard.

\subsection{Architecture}

The current architecture can be represented as:

+-----------------------------+
\begin{table}[h!]
\centering
\resizebox{\textwidth}{!}{%
\begin{tabular}{|l|}
\hline
React Frontend \\ \hline
Dashboard / Explorer / UI \\ \hline
\end{tabular}
}
\end{table}
+--------------+--------------+
\begin{table}[h!]
\centering
\resizebox{\textwidth}{!}{%
\begin{tabular}{|}
\hline
 \\ \hline
HTTP / JSON \\ \hline
\end{tabular}
}
\end{table}
v
+-----------------------------+
\begin{table}[h!]
\centering
\resizebox{\textwidth}{!}{%
\begin{tabular}{|l|}
\hline
Express API Bridge \\ \hline
Routes, parsing, logs, port \\ \hline
scanning, core execution \\ \hline
\end{tabular}
}
\end{table}
+--------------+--------------+
\begin{table}[h!]
\centering
\resizebox{\textwidth}{!}{%
\begin{tabular}{|}
\hline
 \\ \hline
child\_process exec \\ \hline
\end{tabular}
}
\end{table}
v
+-----------------------------+
\begin{table}[h!]
\centering
\resizebox{\textwidth}{!}{%
\begin{tabular}{|l|}
\hline
C Core Binaries \\ \hline
viewer.exe, node\_app.exe, \\ \hline
blockchain.exe \\ \hline
\end{tabular}
}
\end{table}
+--------------+--------------+
\begin{table}[h!]
\centering
\resizebox{\textwidth}{!}{%
\begin{tabular}{|}
\hline
 \\ \hline
\end{tabular}
}
\end{table}
+-----------------------+-----------------------+
\begin{table}[h!]
\centering
\resizebox{\textwidth}{!}{%
\begin{tabular}{|l|}
\hline
 \\ \hline
\end{tabular}
}
\end{table}
v                                               v
+-------------------+                         +---------------------+
\begin{table}[h!]
\centering
\resizebox{\textwidth}{!}{%
\begin{tabular}{|l|l|l|}
\hline
blockchain.dat &  & offchain/records \\ \hline
append-only chain &  & encrypted files \\ \hline
\end{tabular}
}
\end{table}
+-------------------+                         +---------------------+

The distributed demo network can be represented as:

+-------------+        TCP         +-------------+
\begin{table}[h!]
\centering
\resizebox{\textwidth}{!}{%
\begin{tabular}{|l|l|l|}
\hline
Node 8001 & <---------------> & Node 8002 \\ \hline
Validator &  & Validator \\ \hline
\end{tabular}
}
\end{table}
+------+------+                   +------+------+
^                                 ^
\begin{table}[h!]
\centering
\resizebox{\textwidth}{!}{%
\begin{tabular}{|l|}
\hline
 \\ \hline
TCP & TCP \\ \hline
\end{tabular}
}
\end{table}
v                                 v
+-------------+        TCP         +-------------+
\begin{table}[h!]
\centering
\resizebox{\textwidth}{!}{%
\begin{tabular}{|l|l|l|}
\hline
Node 8003 & <---------------> & Node 8004 \\ \hline
Validator &  & Validator \\ \hline
\end{tabular}
}
\end{table}
+-------------+                   +-------------+

Each node has its own:

\begin{itemize}[leftmargin=1.8em]
\item node\_app.exe binary
\item keys directory
\item data/blockchain.dat
\item data/network.log
\item offchain/records directory
\end{itemize}

The API currently reads the live demo node path first, especially demo\_instances/node1, and falls back to backend/core/data in some situations.

\subsection{Implementation Details}

3.3.1 C Blockchain Core

The C core is the foundation of the system.

3.3.1.1 Main Data Structures

The Transaction structure stores:

\begin{itemize}[leftmargin=1.8em]
\item patient\_id
\item doctor\_id
\item data\_hash
\item data\_pointer
\item timestamp
\end{itemize}

The Block structure stores:

\begin{itemize}[leftmargin=1.8em]
\item index
\item timestamp
\item previous\_hash
\item block\_hash
\item merkle\_root
\item validator\_signature
\item validator\_port
\item transactions
\item transaction\_count
\end{itemize}

These structures are fixed-size and simple. This is useful for deterministic file storage and easier serialization, although it also creates limits, such as fixed maximum transaction count and fixed string lengths.

3.3.1.2 Hashing

The project includes a standalone SHA-256 implementation in C. It is used for:

\begin{itemize}[leftmargin=1.8em]
\item hashing text data
\item hashing file contents
\item computing transaction hashes
\item computing Merkle roots
\item computing block hashes
\end{itemize}

The block hash is computed from:

\begin{itemize}[leftmargin=1.8em]
\item block index
\item timestamp
\item previous hash
\item Merkle root
\end{itemize}

The Merkle root is computed from transaction fields. This means that if a transaction changes, the Merkle root changes, and therefore the block hash changes.

3.3.1.3 Genesis Block

The chain begins with a genesis block. The current implementation uses a fixed timestamp and a universal genesis validator port of 8001. The genesis transaction uses patient\_id = GENESIS and doctor\_id = NETWORK. This creates a predictable starting point for all nodes.

3.3.1.4 Persistence

Blocks are appended to a binary file:

data/blockchain.dat

The add\_block function writes the block, flushes the file, and uses an explicit disk sync:

\begin{itemize}[leftmargin=1.8em]
\item \_commit on Windows
\item fsync on Unix-like systems
\end{itemize}

This is good because it reduces the risk of data remaining only in memory after a write.

3.3.1.5 Verification

The core verifies:

\begin{itemize}[leftmargin=1.8em]
\item block hash correctness
\item expected demo signature string validity in the current non-production signature implementation
\item previous-hash linkage
\item duplicate transaction hashes inside a block
\item duplicate transaction hashes against the existing chain
\end{itemize}

Full chain verification reads blocks sequentially and checks that each block points to the previous stored hash.

3.3.1.6 Duplicate Detection

The project includes a transaction hash index implemented as a hash map with chained linked lists. The index stores the transaction data\_hash together with the block index and transaction index. This improves duplicate lookup compared to scanning the whole chain every time and also supports proof lookup for VERIFY\_FILE. Recent corrections clear the index when the blockchain file path changes, which matters because demo nodes may use different ledger files. In production this index should be persistent and stored in a database, but clients must still verify returned proofs rather than trusting the index server blindly.

3.3.1.7 Viewer Tool

The viewer binary provides commands such as:

\begin{itemize}[leftmargin=1.8em]
\item HEIGHT
\item LAST
\item ALL
\item VERIFY
\item VERIFY\_FILE
\item ADD
\end{itemize}

VERIFY\_FILE now returns a proof-backed response for a matching file. It recomputes the uploaded file hash, looks up the transaction location through the in-memory hash index, loads the relevant block, generates a Merkle proof, verifies the proof against the block Merkle root inside the C core, and prints header, transaction, and proof fields for the API layer to parse.
\begin{itemize}[leftmargin=1.8em]
\item PRINT <index>
\end{itemize}

The Express API depends heavily on this viewer output. It parses the text format and converts it to JSON for the frontend.

3.3.1.8 Key Generator and Key Management Architecture

The current implementation splits key creation and management into two paths:

1. Validator Keys: Created via the dedicated helper utility generate\_keys.c. Running this program generates mock public and private key files on disk (e.g., keys/\{port\}\_private.pem and keys/\{port\}\_public.pem). In a production blockchain network, validators would generate real RSA-4096 or ECDSA private keys locally, submitting the corresponding public keys to the consortium authority for inclusion in the authorized validator set.
2. Patient and Doctor Keys: In this prototype version, no standalone cryptographic key files are generated for end-users. Instead, their unique database identifiers (such as PAT\_001 or DOC\_001) act directly as their public keys. The backend API maps the user's authenticated session to these identifiers, which are then used as parameters for the C core's XOR-based mock envelope encryption. In a production environment, this is replaced by client-side keypairs stored in device hardware enclaves.

3.3.2 Network and Consensus Layer

The network layer is implemented in C using TCP sockets. It supports both Windows and Unix-like environments through conditional compilation.

3.3.2.1 Peer Connections

Each node starts a server socket and accepts incoming peer connections. It can also connect to peer ports passed through the command line. Nodes identify themselves using MY\_PORT messages.

3.3.2.2 Synchronization

The synchronization flow includes:

1. A node requests or receives CHAIN\_HEIGHT.
2. If a peer has a higher height, the local node requests missing blocks using GET\_BLOCK:<index>.
3. The peer sends SYNC\_BLOCK messages.
4. The receiving node deserializes, verifies, and appends missing blocks.

This design allows a lagging node to catch up from peers.

3.3.2.3 Three-Phase-Commit-Inspired Flow

The most important current network flow is:

1. Proposer creates a block.
2. Proposer broadcasts CAN\_COMMIT:<serialized block>.
3. Peers verify the block.
4. Peers send VOTE\_YES:<port>.
5. Proposer waits for enough unique votes.
6. Proposer sends PRE\_COMMIT.
7. Peers reply ACK\_PRE:<port>.
8. Proposer waits for enough ACK messages.
9. Proposer appends the block and broadcasts DO\_COMMIT.
10. Peers append their pending block when DO\_COMMIT arrives.

ASCII sequence:

Proposer                 Peer 1                 Peer 2
\begin{table}[h!]
\centering
\resizebox{\textwidth}{!}{%
\begin{tabular}{|l|l|}
\hline
 &  \\ \hline
--- CAN\_COMMIT -------> &  \\ \hline
--- CAN\_COMMIT -----------------------------> \\ \hline
 &  \\ \hline
<-- VOTE\_YES ---------- &  \\ \hline
<-------------------------------- VOTE\_YES --- \\ \hline
 &  \\ \hline
--- PRE\_COMMIT -------> &  \\ \hline
--- PRE\_COMMIT -----------------------------> \\ \hline
 &  \\ \hline
<-- ACK\_PRE ----------- &  \\ \hline
<-------------------------------- ACK\_PRE ---- \\ \hline
 &  \\ \hline
--- DO\_COMMIT --------> &  \\ \hline
--- DO\_COMMIT ------------------------------> \\ \hline
 &  \\ \hline
\end{tabular}
}
\end{table}
append locally          append locally          append locally

3.3.2.4 Current Consensus Strength

This is a useful prototype consensus workflow. It shows that peers do not blindly accept any block. They validate and vote. It also avoids accepting duplicate votes through a vote mask.

3.3.2.5 Current Consensus Limitations

The current consensus is not production Byzantine fault tolerant. It assumes mostly honest nodes. It does not solve all cases of network partition, equivocation, malicious proposer behavior, view change, validator slashing, or formal leader election. The majority threshold is currently coded around a small demo network. A production system would need a configurable validator set and formal quorum rules.

3.3.3 Backend API Bridge

The backend API is a Node.js/Express service.

3.3.3.1 Purpose

The API exists because the frontend cannot directly read binary blockchain files or execute C programs safely. The API acts as a bridge between browser and core.

3.3.3.2 Main Routes

The blockchain routes include GET /api/height, GET /api/blocks, GET /api/block/:index, and GET /api/verify.

The network routes include GET /api/network/health and GET /api/network/logs.

The record route includes:

\begin{itemize}[leftmargin=1.8em]
\item POST /api/records/upload
\end{itemize}

3.3.3.3 Core Bridge

The API runs the C viewer executable through child\_process.exec. It first attempts to run commands from a live demo node directory, especially demo\_instances/node1. It then parses viewer output.

This approach is simple and effective for a demo. It also keeps the C core separate from the JavaScript server. However, it is not ideal for production because child process execution is slower and riskier than a structured native interface, RPC service, or database-backed bridge.

3.3.3.4 Parser

The parser reads lines such as:

INDEX|...
TIMESTAMP|...
PREV\_HASH|...
BLOCK\_HASH|...
MERKLE|...
VALIDATOR|...
SIGNATURE|...
TX\_COUNT|...
TX|...
---END\_BLOCK---

It converts them into JSON blocks with nested transactions.

3.3.3.5 Network Health

The network health controller scans ports 8001 through 8010 on localhost. Any port accepting a TCP connection is treated as an online node. This is practical for the local demo and dashboard. For production, node discovery should come from a signed registry or service discovery layer rather than port scanning.

3.3.3.6 Logs

The API reads:

\begin{itemize}[leftmargin=1.8em]
\item backend API server logs
\item demo\_instances/node*/data/network.log
\end{itemize}

It aggregates these logs for the frontend System Logs page.

3.3.3.7 Authentication

The backend and frontend authentication models have been unified around a robust, zero-dependency local mock authentication system. The backend middleware validates the X-Mock-User-ID header against an internal mock user registry. This removes the integration gap that previously existed when mixing local frontend state with external Supabase dependencies, ensuring a fully self-contained local demonstration.

3.3.4 Frontend Dashboard

The frontend is a React TypeScript dashboard.

3.3.4.1 Main User Roles

The application distinguishes between patient, doctor, and admin roles.

Each role sees different sidebar links and allowed pages.

3.3.4.2 Dashboard

The dashboard displays the chain height, network status, active node count, and recent blocks.

This gives a quick operational overview.

3.3.4.3 Block Explorer

The block explorer displays blocks, hashes, previous hashes, Merkle roots, validator ports, transaction count, and metadata. It also includes a "Verify Full Chain" action that calls the backend verification route.

3.3.4.4 Network Health

The network page displays online validator nodes, ports, type, status, and latency. It is useful for showing that the blockchain is not only a data structure but a distributed network.

3.3.4.5 System Logs

The system logs page merges API logs and node logs. This supports audit and debugging.

3.3.4.6 Upload Record

The upload page lets a doctor or admin enter:

\begin{itemize}[leftmargin=1.8em]
\item patient ID
\item doctor ID
\item encryption password
\item medical file
\end{itemize}

It sends the file through FormData to the /api/records/upload endpoint. The backend calls the C core ENCRYPT command to securely encrypt the record under the doctor's and patient's key envelopes, saving the resulting .enc file in the demo node's offchain/records directory. It then calls the C core's ADD command with the file path, data pointer, patient\_id, doctor\_id, and validator port. On success, the frontend displays the committed block index, SHA-256 hash, data pointer, patient ID, and doctor ID returned by the core. Duplicate file submissions are detected and rejected with a 409 response.

3.3.4.7 My Records

The records page reads blockchain blocks and extracts transactions. If the user is a patient, it filters records by patient\_id. Admins can see all records.

3.3.4.8 Verify Record

The verify page uploads the file to /api/records/verify. The backend passes the file to the C viewer VERIFY\_FILE command. The C core recomputes the SHA-256 hash of the uploaded file, uses the transaction hash index to locate the block index and transaction index, and generates a Merkle proof for the matching transaction. If a matching data\_hash is found, the response includes the block index, transaction index, record hash, data pointer, patient ID, doctor ID, transaction metadata, block header fields, Merkle proof leaf, proof length, computed root, and each sibling proof step with left/right direction.

The frontend no longer treats the backend's MATCH response as sufficient by itself. It recomputes SHA-256 over the uploaded file in the browser, reconstructs the canonical transaction leaf from returned metadata, applies the Merkle proof step by step, compares the computed root with the block header Merkle root, and recomputes the block hash from the returned header fields. The page displays success only if the local proof check succeeds. If no match is found, the computed hash is shown alongside a clear failure message.

3.3.4.9 Activity Logs

The activity logs page converts block and transaction history into audit events such as BLOCK\_SIGNED and RECORD\_UPLOAD.

3.3.4.10 Settings

The settings page is mostly a static UI placeholder for profile, security, notifications, storage configuration, and public key display.

3.3.5 Data Model and Storage Design

3.3.5.1 On-Chain Data

The chain stores:

\begin{itemize}[leftmargin=1.8em]
\item block index
\item block timestamp
\item previous block hash
\item block hash
\item Merkle root
\item validator signature
\item validator port
\item transaction array
\end{itemize}

Each transaction stores:

\begin{itemize}[leftmargin=1.8em]
\item patient\_id
\item doctor\_id
\item data\_hash
\item data\_pointer
\item timestamp
\end{itemize}

3.3.5.2 Off-Chain Data

Encrypted medical files are stored in offchain/records folders. The data\_pointer field points to these files. The blockchain does not need the file contents to preserve integrity; it only needs the file hash.

3.3.5.3 Why Hybrid Storage Is Important

If raw medical files were stored directly on-chain:

\begin{itemize}[leftmargin=1.8em]
\item every node would store patient data
\item deletion would become extremely difficult
\item privacy exposure would increase
\item chain size would grow rapidly
\item compliance concerns would become serious
\end{itemize}

With hybrid storage:

\begin{itemize}[leftmargin=1.8em]
\item the chain remains lightweight
\item private files remain outside the ledger
\item file integrity can still be verified
\item deletion or access revocation is more realistic
\end{itemize}

3.3.5.4 Merkle Root

The Merkle root summarizes all transactions in a block. If a transaction changes, the Merkle root changes. Since the block hash includes the Merkle root, transaction tampering changes the block hash and breaks verification.

3.3.5.5 Merkle Proof Verification

The implementation now supports proof-backed record verification. A Merkle proof contains the transaction leaf hash, each sibling hash needed to climb from the leaf to the root, and the left/right position of every sibling. This allows a lightweight client to verify record inclusion using the uploaded file, transaction metadata, proof path, and block header instead of downloading the full block.

The current frontend performs this proof reconstruction in the browser. This provides inclusion verification, but production SPV security also requires validation that the returned block header is part of the authorized chain.

3.3.6 Transaction Lifecycle

The ideal transaction lifecycle is:

1. Doctor selects a medical file.
2. System encrypts or receives an encrypted file.
3. System computes SHA-256 hash of the encrypted file.
4. System creates a transaction containing patient\_id, doctor\_id, data\_hash, pointer, and timestamp.
5. Node checks whether data\_hash already exists.
6. Node creates a block.
7. Node computes Merkle root.
8. Node computes block hash.
9. Validator signs the block hash.
10. Node broadcasts CAN\_COMMIT.
11. Peers verify block.
12. Peers vote yes if valid.
13. Proposer sends PRE\_COMMIT.
14. Peers acknowledge.
15. Proposer and peers commit the block.
16. Frontend dashboard displays the new block and record.
17. Later, a verifier recomputes the file hash and compares it to the stored hash.

Current implementation status:

\begin{itemize}[leftmargin=1.8em]
\item Steps 3 through 15 exist in the C core/node flow for command-line record addition.
\item Dashboard display exists.
\item Upload from frontend to backend is now fully integrated: the API controller receives the file, calls the C core ENCRYPT command to securely encrypt the record under patient and doctor envelopes in the offchain/records directory, and then calls the ADD command to commit the transaction to the blockchain. The committed block data is then returned to the frontend UI.
\item Frontend verification is now integrated: the file is uploaded to the backend, the C viewer VERIFY\_FILE command recomputes the hash, locates the transaction through the hash index, returns a Merkle proof and block header, and the browser verifies the proof locally before displaying success.
\end{itemize}

\subsection{Advantages over Existing Systems}

The proposed Permissioned Blockchain-Based Medical Record Integrity System intentionally strips away the unnecessary computational bloat of generalized platforms, synthesizing the lessons learned from the aforementioned literature to provide a pragmatic, full-stack operational solution.

\subsubsection{Eradicating the Single Point of Failure and Insider Threats}
Rather than relying on probabilistic AI anomaly detection to catch insider threats, the proposed system resolves the centralized database crisis through deterministic cryptographic immutability.
\textbf{The Solution:} The project implements a custom-coded, three-phase-commit-inspired (3PC) consensus protocol (CAN\_COMMIT, VOTE\_YES, PRE\_COMMIT, ACK\_PRE, DO\_COMMIT) over raw TCP sockets. This ensures synchronous agreement, meaning a single malicious database administrator cannot unilaterally alter a patient's history. This Proof of Authority (PoA) consensus style—verified as highly efficient for healthcare environments by recent frameworks like BlockCare (2024)—entirely eliminates the single point of failure without the energy waste of PoW.

\subsubsection{Bypassing EVM Constraints with a Native C-Core}
Unlike MedRec or Ancile which suffer from Ethereum gas fees and 256-bit word limits, this project completely circumvents public network bottlenecks.
\textbf{The Solution:} The architecture executes a low-level, native C blockchain core. It enforces deterministic memory allocation by defining rigid, fixed-width C struct representations for blocks and transactions, storing them directly to disk via blockchain.dat with explicit fsync operations. This guarantees predictable, zero-cost computational performance, allowing the network to scale rapidly.

\subsubsection{Separation of Concerns: Decoupling Logic from the Ledger}
To avoid the computational bloat seen in FHIRChain and MedChain, the proposed system aggressively minimizes its on-chain footprint.
\textbf{The Solution:} The blockchain transaction struct only stores foundational metadata (patient\_id, doctor\_id, data\_hash, data\_pointer, timestamp). The heavy lifting of role-based access control (RBAC), user session management, and interface rendering is entirely decoupled from the C core and handled by the Node.js Express API and Vite React TypeScript frontend. By acting as a mathematically pristine engine of integrity rather than a bloated application server, the system achieves horizontal scalability.

\subsubsection{Deterministic Local Off-Chain Storage (Replacing IPFS Risks)}
To resolve the latency and data-loss risks of IPFS swarms (seen in BiiMED and EHRChain) and the centralization of cloud providers (seen in MeDShare), the system implements a natively managed hybrid storage architecture.
\textbf{The Solution:} When a record is uploaded, the Node.js API triggers the C core's ENCRYPT command, which executes mock envelope encryption (slated for upgrade to AES-256-GCM). Crucially, the resulting .enc file is stored in a localized offchain/records directory physically located on the node's disk. This guarantees absolute high availability, deterministic retrieval speeds, and strict HIPAA compliance, as the file resides on controlled hospital infrastructure rather than an ephemeral public swarm.

\subsubsection{Deterministic Cryptographic Duplicate Prevention}
The system replaces the computationally heavy semantic matching of traditional MDM frameworks highlighted by Tayal and Adapa with mathematical precision.
\textbf{The Solution:} The architecture implements a transaction hash index engineered with a pure-C Bloom filter utilizing three independent DJB2-based hash functions. This provides a constant-time O(1) pre-filtering cache. If a duplicate file hash is detected, the API instantly rejects the upload with an HTTP 409 Conflict response before it hits the consensus network. To scale this determinism to over 1 billion records without exhausting physical RAM, the project outlines a LevelDB Log-Structured Merge (LSM) tree integration, bounding memory usage to ~2 GB.

\subsubsection{Zero-Trust Browser-Side Merkle Verification}
A critical vulnerability in web-based blockchain dashboards is the light-client trust boundary—if the backend API is compromised, it can falsely validate tampered records.
\textbf{The Solution:} The proposed project implements zero-trust Merkle proof verification directly within the browser. When verifying a record, the C core generates a full Merkle proof containing sibling hashes and directional paths. The React frontend independently recomputes the SHA-256 hash of the uploaded file via Web Crypto APIs, reconstructs the canonical leaf, applies the proof to climb the tree, and mathematically verifies the root against the block header. This ensures end-to-end integrity, proving that even a compromised middleware layer cannot spoof mathematical truth.


3.4.7 System Strengths

\begin{itemize}[leftmargin=1.8em]
\item Clear Separation of Concerns
\end{itemize}

The project separates:

\begin{itemize}[leftmargin=1.8em]
\item blockchain core
\item API bridge
\item frontend dashboard
\item documentation
\item scripts
\item demo instances
\end{itemize}

This makes the codebase easier to understand.

\begin{itemize}[leftmargin=1.8em]
\item Strong Educational Value
\end{itemize}

The project exposes many important concepts:

\begin{itemize}[leftmargin=1.8em]
\item hashing
\item block chaining
\item Merkle roots
\item block verification
\item peer networking
\item synchronization
\item role-based UI
\item API bridging
\item off-chain storage
\item audit logs
\end{itemize}

\begin{itemize}[leftmargin=1.8em]
\item Practical Healthcare-Oriented Design
\end{itemize}

The hybrid model is correct for healthcare. Storing hashes on-chain and files off-chain is much better than storing raw records on-chain.

\begin{itemize}[leftmargin=1.8em]
\item Full Stack Demonstrability
\end{itemize}

The project can be shown visually through the dashboard. This makes it stronger than a command-line-only thesis project.

\begin{itemize}[leftmargin=1.8em]
\item Low Dependency Core
\end{itemize}

The current C core avoids heavy external cryptographic dependencies in demo mode. This makes compilation easier on Windows. It also makes the internal flow easier to inspect.

\begin{itemize}[leftmargin=1.8em]
\item Explicit Disk Sync
\end{itemize}

Using flush and disk sync after appending a block is a good reliability practice.

\begin{itemize}[leftmargin=1.8em]
\item Duplicate Record Handling
\end{itemize}

The transaction hash index is a meaningful improvement. It turns duplicate detection from a purely linear lookup into indexed lookup.

\begin{itemize}[leftmargin=1.8em]
\item Multi-Node Demo Structure
\end{itemize}

The demo\_instances folder supports multiple local nodes, each with its own data, keys, logs, and records. This is useful for demonstration.

\begin{itemize}[leftmargin=1.8em]
\item Dashboard Breadth
\end{itemize}

The frontend covers many expected operational views:

\begin{itemize}[leftmargin=1.8em]
\item dashboard
\item explorer
\item network
\item logs
\item upload
\item records
\item verification
\item settings
\end{itemize}

This gives the project product depth.


\newpage
\section{SECURITY MODEL AND EVALUATION}

\subsection{Security Model}

4.1.1 Security Goals

The project targets:

\begin{itemize}[leftmargin=1.8em]
\item integrity
\item tamper evidence
\item validator accountability
\item duplicate detection
\item auditability
\item privacy by not storing raw medical data on-chain
\end{itemize}

4.1.2 Tamper Evidence

Tamper evidence is provided by:

\begin{itemize}[leftmargin=1.8em]
\item SHA-256 transaction hashes
\item Merkle roots
\item previous\_hash linkage
\item full chain verification
\end{itemize}

Changing one stored block changes its hash. The next block still points to the old hash, so verification fails.

4.1.3 Duplicate Detection

The system indexes transaction hashes. If the same file hash already exists, the node can reject the record. This helps prevent replay of identical medical documents.

4.1.4 Validator Identity

Each block stores validator\_port and validator\_signature. In a production system, validator identity should be tied to a real public key, certificate, and institution identity. In the current demo, signature verification is simplified.

4.1.5 Current Signature Limitation

The current signature implementation is not cryptographically secure. sign\_data creates a string beginning with SIG\_V1\_ and includes the block hash and key path. verify\_signature now reconstructs the expected demo signature for the supplied block hash and corresponding key path rather than accepting only a prefix, but this is still a demonstration mechanism rather than real public-key cryptography. It is acceptable only for demonstrating the block signing workflow. It is not sufficient for real security.

A production version should use:

\begin{itemize}[leftmargin=1.8em]
\item Ed25519, ECDSA, or RSA-PSS signatures
\item real public/private key verification
\item certificate management
\item signed block headers for light-client validation
\item key rotation
\item revocation list
\item hardware-backed private keys for validators
\end{itemize}

4.1.6 Merkle Proof and Light-Client Trust Boundary

The current verifier can prove that a returned transaction belongs to a returned block header. The browser recomputes the file hash, reconstructs the transaction leaf, applies the left/right sibling proof, and checks the computed root against the returned Merkle root. This reduces the need to download full blocks for record inclusion checks.

However, Merkle proof verification alone does not prove that the block header belongs to the real authorized chain. A production light client must also validate the block header through trusted validator public keys, signed header chains, previous\_hash linkage, validator-set rules, and/or trusted checkpoints. Without this additional header validation, a malicious server could return a fake header and a matching fake proof.

4.1.7 Network Security

Current node communication is plain TCP. It does not use TLS. This means traffic can be observed or modified by an attacker on the network. The block verification layer can detect some invalid data, but confidentiality and authenticity of transport are not guaranteed.

Production should use:

\begin{itemize}[leftmargin=1.8em]
\item TLS or mutual TLS
\item node certificates
\item signed peer registry
\item replay protection
\item message authentication
\item rate limiting
\end{itemize}

4.1.7 Backend Security

The API uses child\_process execution to call C binaries. This is acceptable for a controlled demo but requires hardening. Commands must be strictly controlled to avoid injection risk. In the current code, commands are mostly predefined, but future features should avoid passing raw user input into shell commands.

4.1.8 Frontend Security

The system currently uses local mock authentication across both frontend and backend. This is suitable for demonstration and local testing but not for real deployment. Production must integrate a real Identity Provider (IdP) consistently across the stack.

\subsection{Testing and Verification Status}

During the latest review, the verification commands passed successfully, including frontend linting and building, backend API JavaScript syntax checks, and core compilation via MinGW. Additionally, the core viewer correctly returns EMPTY when no central core ledger exists.

The frontend build produced only a Vite bundle-size warning, not a failure.

Important interpretation:

The core viewer returns EMPTY when no data/blockchain.dat exists in backend/core. Demo instances still contain their own chain files. Therefore, an empty central core chain is not the same as a corrupted demo chain.

Current testing gaps include the absence of API unit tests, frontend component tests, automated multi-node consensus tests, upload integration tests, and long-running network stability tests. Cryptographic signature tests are also omitted because the current signatures are mock implementations, and there is no fuzz testing for the serializer/parser.


\newpage
\section{DISCUSSION AND LIMITATIONS}

\subsection{Current Weaknesses and Gaps}

5.1.1 Mock Signatures

The largest security weakness is that the signature system is not real cryptography. It demonstrates the shape of signing but does not protect against forgery.

5.1.2 Auth Model Is Demo-Level

The backend and frontend authentication have been unified around a local mock registry via the X-Mock-User-ID header. This provides a clean, zero-dependency demonstration experience without requiring external services like Supabase. However, for a production system, this must be replaced with a single, secure Identity Provider (IdP) with proper token validation, session management, and role enforcement backed by a real identity store.

5.1.3 Upload Flow Is Integrated at Demo Level

The upload controller is now fully integrated with the C core encryption commands. When a record is uploaded, the Express API calls the C core's ENCRYPT command (which implements Mock Envelope Encryption). This derives a temporary file key, encrypts it under both the patient's ID and doctor's ID, encrypts the file payload, and stores it in the demo node's offchain/records directory as a secure .enc file. The API then calls the ADD command to commit the transaction to blockchain.dat. The remaining limitation is that the cryptographic operations use a mock XOR cipher (for zero-dependency compilation on Windows) rather than standard production-grade AES-256-GCM and RSA-4096.

5.1.4 Verification UI Is Now Integrated

The Verify Record page no longer simulates verification. It uploads the file to /api/records/verify. The backend passes the file to the C viewer VERIFY\_FILE command, which recomputes the SHA-256 hash, uses the hash index to locate the transaction, and returns a proof-backed response. A MATCH response returns the block index, transaction index, hash, data pointer, patient ID, doctor ID, transaction metadata, block header fields, and Merkle proof steps. A NOT\_FOUND response returns the computed hash. The frontend recomputes the file hash and verifies the returned Merkle proof locally before displaying success.

5.1.5 Plain TCP Networking

Node communication is not encrypted. Production healthcare systems need TLS or mutual TLS.

5.1.6 Consensus Is Demo-Level

The three-phase commit inspired logic is useful but not a complete consensus protocol. It lacks formal handling for malicious validators, network partitions, leader changes, conflicting proposals, and dynamic validator sets.

5.1.7 Binary File Storage Has Portability Risks

The ledger writes C structs directly to disk. This is simple, but it can be sensitive to compiler, platform, architecture, struct layout, and time\_t size differences. A production ledger should use a canonical binary or JSON/CBOR/Protobuf format.

5.1.8 Limited Automated Tests

The repository contains test sources, but the frontend and backend do not have meaningful automated test suites. backend/api npm test is a placeholder.

5.1.9 Logs and Demo Data Are Tracked

Demo blockchain files, node logs, and binaries appear in the repository. For production-quality repository hygiene, generated artifacts should usually be excluded or separated from source.

5.1.10 Mock Hashing and Cipher Limitations

While the end-to-end encryption pipeline is now fully integrated into both web uploads and downloads via the C core's ENCRYPT/DECRYPT commands, it uses simulated cryptography (XOR-based mock ciphers) rather than standard OpenSSL cryptosystems. This is suitable for demonstrating the structural envelope flow, but does not provide cryptographically secure data protection.

5.1.11 No Real Patient Consent Model

A medical system should record patient consent, access grants, revocations, and audit trails. The current project focuses on integrity, not consent governance.

5.1.12 No Production Database Schema

A robust identity and access management system must be implemented. The chain and frontend are not fully integrated with a stable database schema for users, permissions, records, and audit metadata.

\subsection{Compliance and Ethical Considerations}

5.2.1 Privacy

The system correctly avoids putting raw medical files on-chain. This is the most important privacy-preserving decision in the architecture.

5.2.2 Data Minimization

Only hashes, pointers, and limited metadata should be stored on-chain. In future versions, patient\_id and doctor\_id should be pseudonymized or tokenized more strongly.

5.2.3 Right to Erasure

Blockchain data is immutable, so deleting on-chain data is difficult. The hybrid model helps because the medical file can be deleted off-chain while the chain retains only a hash. This is often called crypto-shredding when encryption keys or encrypted files are destroyed.

5.2.4 Consent

The current implementation does not fully model patient consent. A production medical system should record:

\begin{itemize}[leftmargin=1.8em]
\item consent granted
\item consent revoked
\item consent scope
\item who accessed a record
\item why access occurred
\item emergency override events
\end{itemize}

5.2.5 Auditability

The ledger and logs support auditability, but production audit logs need stronger guarantees. API logs should be append-only, protected from modification, and possibly anchored into the blockchain itself.

5.2.6 Ethical Use

A system like this should never expose raw medical data unnecessarily. It should also avoid giving a false sense of security. Blockchain protects integrity, but it does not automatically solve access control, confidentiality, consent, or clinical correctness.


\newpage
\section{FUTURE WORK AND PRODUCTION ROADMAP}

\subsection{Future Improvements}

6.1.1 Replace Mock Signatures with Real Cryptography

This is the highest priority.

Recommended cryptographic options include Ed25519 for fast modern signatures, ECDSA P-256 for standards compatibility, or RSA-PSS if RSA is required.

The implementation should verify that a block was signed by the private key corresponding to the validator public key.

6.1.2 Use Canonical Serialization

Instead of writing C structs directly to disk, the system should use a stable format such as Protocol Buffers, CBOR, FlatBuffers, or a custom fixed-width binary format.

This would improve cross-platform compatibility.

6.1.3 Transition to Standard Production Cryptography and Secure Key Management

While the C core and Express API now implement a fully integrated upload and download flow featuring Mock Envelope Encryption, a production system must transition from simulated (XOR-based) cryptography to standard production algorithms:

1. Replace mock XOR with real AES-256-GCM symmetric encryption for off-chain record payloads.
2. Replace mock RSA with real RSA-4096 or ECIES (Elliptic Curve Integrated Encryption Scheme) for envelope key wrapping.
3. Integrate a standard Key Management Service (KMS) or HSM (Hardware Security Module) to handle validator and administrative keys securely, rather than relying on key files in local directories.
4. Establish Non-Custodial Key Management for Patients and Doctors: Storing patient/doctor private keys in a centralized server database is unsafe, and expecting patients to memorize raw 256-bit hexadecimal keys is impractical. Future production systems must employ:
\begin{itemize}[leftmargin=1.8em]
\item BIP-39 Mnemonic Seed Phrases: Patients are given 12 or 24 words from which their private key is derived on-the-fly client-side, removing any need to memorize raw keys.
\item Hardware enclaves: Private keys are locked inside device-level Secure Enclaves (e.g., Apple Secure Enclave, Android Keystore, or TPMs), accessible only via biometric Touch ID/Face ID verification, ensuring keys never leave the device.
\item Threshold social recovery: Utilizing Shamir's Secret Sharing to split keys into parts across the cloud, local device, and trusted consortium guardians so key recovery is resilient and decentralized.
\end{itemize}

6.1.4 Implement Real Verification

The Verify Record page now:

\begin{itemize}[leftmargin=1.8em]
\item uploads the file to the backend for lookup and proof generation,
\item computes SHA-256 in the browser,
\item reconstructs the canonical transaction leaf,
\item applies the returned Merkle proof step by step,
\item compares the computed root with the returned block header Merkle root,
\item recomputes the block hash from the returned header fields.
\end{itemize}

Remaining production work is to validate the returned block header against real validator signatures, trusted public keys, previous\_hash linkage, and a checkpoint/header-chain policy.

6.1.5 Unify Authentication

Choose one auth model:

\begin{itemize}[leftmargin=1.8em]
\item Production Identity Provider (IdP) integration end to end, or
\item local mock auth only for demo, or
\item JWT issued by backend
\end{itemize}

The frontend and backend must agree on token format and role source.

6.1.6 Add TLS / Mutual TLS

Node-to-node communication should be encrypted and authenticated. In a hospital consortium, each validator should have a certificate.

6.1.7 Add Validator Governance

The system needs:

\begin{itemize}[leftmargin=1.8em]
\item validator registry
\item onboarding process
\item removal process
\item key rotation
\item revocation list
\item quorum configuration
\end{itemize}

6.1.8 Improve Consensus

Options:

\begin{itemize}[leftmargin=1.8em]
\item formal Raft-style leader election for crash fault tolerance
\item PBFT/Tendermint-style consensus for Byzantine settings
\item configurable quorum rules
\item timeout and recovery from failed proposals
\item protection against double proposals
\end{itemize}

6.1.9 Implement LSM-Tree Disk-Backed Indexing (LevelDB Integration)

To scale the ledger to 1 billion records without exhausting physical RAM (which would otherwise require over 184 GB of memory for pure in-memory indexes), the current in-memory separate chaining hash maps should be migrated to LevelDB or RocksDB. The proposed Log-Structured Merge (LSM) tree key-value database schema would support:
1. Transaction Hash Index: Key 'tx:<data\_hash>' -> Value '<block\_index>:<tx\_index>'
2. Patient History Index: Key 'pat:<patient\_id>:<block\_index>:<tx\_index>' -> Value ''
3. Doctor Audit Index: Key 'doc:<doctor\_id>:<block\_index>:<tx\_index>' -> Value ''

LevelDB's MemTable write buffering, sorted SSTables, and built-in Bloom filter capabilities would cap the memory footprint to ~2 GB RAM regardless of ledger size, maintaining O(log N) SSD lookup speeds.

6.1.10 Add Automated Tests

Recommended tests:

\begin{itemize}[leftmargin=1.8em]
\item hash function tests against known SHA-256 vectors
\item block hash reproducibility tests
\item duplicate detection tests
\item chain verification tests
\item serializer/deserializer tests
\item API route tests
\item frontend page rendering tests
\item multi-node sync tests
\item upload integration tests
\end{itemize}

6.1.11 Improve Repository Hygiene

Generated files should be controlled:

\begin{itemize}[leftmargin=1.8em]
\item binaries can be ignored or stored in release artifacts
\item logs should not normally be committed
\item demo data should be explicitly separated from source
\item keys should not be real secrets in the repository
\end{itemize}

6.1.12 Add Monitoring

Production monitoring should include:

\begin{itemize}[leftmargin=1.8em]
\item node online/offline status
\item block height mismatch
\item sync lag
\item failed block verification
\item duplicate record attempts
\item upload failures
\item suspicious access patterns
\end{itemize}

6.1.13 Add Access Control and Consent

Medical data systems need strong user permissions. Future work should include:

\begin{itemize}[leftmargin=1.8em]
\item patient consent grants
\item doctor access requests
\item emergency access mode
\item audit trail for every access
\item role-based and attribute-based access control
\end{itemize}

6.1.14 Add Performance Benchmarking

The thesis should eventually include:

\begin{itemize}[leftmargin=1.8em]
\item multiple benchmark runs
\item average latency
\item standard deviation
\item throughput under different node counts
\item throughput under different transaction counts per block
\item memory usage
\item CPU usage
\item network sync time
\item charts and graphs
\end{itemize}

\subsection{Production Roadmap}

Phase 1 (Stabilize Prototype) focuses on maintaining the integrated upload flow and proof-backed record verification, unifying authentication, adding basic automated tests, and cleaning generated files from the repository.

Phase 2 (Secure Core) involves implementing real signatures, canonical serialization, TLS/mTLS, a signed validator registry, and key rotation and revocation mechanisms.

Phase 3 (Improve Distributed Protocol) introduces a configurable validator set, formal quorums, timeout handling, failed proposer recovery, network partition behavior, and conflict resolution.

Phase 4 (Healthcare Workflow Integration) adds a patient consent model, doctor/patient profile databases, hospital/lab/insurance roles, record access auditing, and integration with existing EHR systems.

Phase 5 (Production Readiness) requires a monitoring dashboard, deployment scripts, containerization, backup and restore procedures, disaster recovery plans, security audits, and comprehensive compliance documentation.


\newpage
\section{CONCLUSION}

\subsection{Summary}

This project is a strong academic and engineering prototype of a permissioned blockchain-based medical record integrity system. Its most important idea is the separation of medical data storage from integrity proof. Sensitive files remain off-chain, while hashes and metadata are stored in a tamper-evident ledger. This design aligns with healthcare privacy needs better than storing raw records directly on-chain.

The current repository demonstrates more than a theoretical blockchain. It includes a C ledger engine, a peer networking layer, a synchronization mechanism, a Node.js API bridge, and a React dashboard. This makes the project suitable for demonstration, explanation, and further development.

The project also honestly reveals the gap between prototype and production. It still needs real cryptographic signatures, trusted light-client header validation, real authentication integration, secure transport, persistent indexing, automated tests, formal validator governance, and stronger consensus rules. These are not minor details; they are essential for real healthcare deployment. However, the current implementation provides a meaningful foundation on which those improvements can be built.

The uniqueness of the work lies in its full-stack nature and its healthcare-specific architecture. It does not merely use blockchain as a buzzword. It uses blockchain for a specific property: tamper evidence. It does not place private data on-chain. It recognizes that permissioned identity is more suitable than public mining for hospitals. It uses a dashboard to make the system understandable to users and examiners.

In final evaluation, the project should be presented as a permissioned medical integrity ledger prototype with a working core, demonstrable dashboard, and clear future path toward production-grade security and compliance.


APPENDIX

23. APPENDIX A: SUGGESTED DIAGRAMS

The following diagrams should be included if this report is converted into a PDF thesis.

\subsection{Overall Architecture Diagram}

React Frontend -> Express API -> C Core Binaries -> Ledger and Off-chain Storage

\subsection{Node Interaction Sequence Diagram}

CAN\_COMMIT -> VOTE\_YES -> PRE\_COMMIT -> ACK\_PRE -> DO\_COMMIT

\subsection{Block Structure Diagram}

+------------------------------------------------+
\begin{table}[h!]
\centering
\resizebox{\textwidth}{!}{%
\begin{tabular}{|l|}
\hline
index \\ \hline
timestamp \\ \hline
previous\_hash \\ \hline
merkle\_root \\ \hline
block\_hash \\ \hline
validator\_port \\ \hline
validator\_signature \\ \hline
transactions[] \\ \hline
\end{tabular}
}
\end{table}
+------------------------------------------------+

\subsection{Transaction Lifecycle Diagram}

File -> Hash -> Transaction -> Block -> Consensus -> Commit -> Verify

\subsection{Threat Model Diagram}

Untrusted client -> API boundary -> Core boundary -> Ledger/off-chain storage

\subsection{Key Flow Diagram}

Validator key file -> sign block hash -> peers verify validator signature

\subsection{Network Topology Diagram}

Local validator nodes on ports 8001-8010 connected through TCP.

24. APPENDIX B: IMPORTANT FILES

Core blockchain:

\begin{itemize}[leftmargin=1.8em]
\item backend/core/src/blockchain/block.h
\item backend/core/src/blockchain/block.c
\item backend/core/src/blockchain/blockchain.c
\item backend/core/src/crypto/hash.c
\item backend/core/src/crypto/signature.c
\item backend/core/src/node\_app.c
\item backend/core/src/viewer.c
\end{itemize}

Network:

\begin{itemize}[leftmargin=1.8em]
\item backend/core/src/network/node.c
\item backend/core/src/network/protocol.c
\item backend/core/src/network/serializer.c
\item backend/core/src/network/proposal.c
\item backend/core/src/network/sync.c
\end{itemize}

API:

\begin{itemize}[leftmargin=1.8em]
\item backend/api/src/index.js
\item backend/api/src/config/core\_bridge.js
\item backend/api/src/controllers/blockController.js
\item backend/api/src/controllers/networkController.js
\item backend/api/src/controllers/recordController.js
\item backend/api/src/utils/parser.js
\item backend/api/src/middleware/auth.js
\end{itemize}

Frontend:

\begin{itemize}[leftmargin=1.8em]
\item frontend/src/App.tsx
\item frontend/src/hooks/useAuth.ts
\item frontend/src/services/api.ts
\item frontend/src/pages/Dashboard.tsx
\item frontend/src/pages/BlockExplorer.tsx
\item frontend/src/pages/NetworkHealth.tsx
\item frontend/src/pages/SystemLogs.tsx
\item frontend/src/pages/UploadRecord.tsx
\item frontend/src/pages/MyRecords.tsx
\item frontend/src/pages/VerifyRecord.tsx
\end{itemize}

26. FORMAL MODEL OF THE LEDGER

\subsection{Block Definition}

Let the blockchain be an ordered sequence:

C = [B0, B1, B2, ..., Bn]

Each block Bi contains:

Bi = (index, timestamp, previous\_hash, merkle\_root, block\_hash, validator\_id, signature, TxList)

For the current implementation, validator\_id is represented by validator\_port.

\subsection{Transaction Definition}

Each transaction T contains:

T = (patient\_id, doctor\_id, data\_hash, data\_pointer, timestamp)

Here:

data\_hash = SHA256(file\_content)

The file\_content should ideally be encrypted medical file content. The blockchain stores the hash, not the raw medical file.

\subsection{Merkle Root}

For each transaction Ti, define:

tx\_hash\_i = SHA256("LEAF" || len(patient\_id) || patient\_id || len(doctor\_id) || doctor\_id || len(data\_hash) || data\_hash || len(data\_pointer) || data\_pointer || timestamp)

The Merkle root is computed by repeatedly hashing pairs:

M = Merkle(tx\_hash\_1, tx\_hash\_2, ..., tx\_hash\_k)

Internal nodes are computed with an explicit node domain:

parent = SHA256("NODE" || left\_hash || right\_hash)

If there is only one transaction, the Merkle root is the transaction hash. If there are multiple transactions, the root summarizes the set. If a level has an odd number of nodes, the final node is duplicated for that level. If any transaction field changes, the root changes.

The implementation can also generate a Merkle proof for a transaction index. The proof contains each sibling hash and whether that sibling is on the left or right side. Verification starts from the transaction leaf and applies each proof step until it reconstructs the root. A reordered proof, wrong sibling hash, or wrong left/right direction should fail.

\subsection{Block Hash}

The block hash is:

block\_hash = SHA256("HEADER" || index || timestamp || previous\_hash || merkle\_root)

This means block integrity depends on:

\begin{itemize}[leftmargin=1.8em]
\item block position
\item block creation time
\item previous block hash
\item transaction summary
\end{itemize}

\subsection{Chain Validity}

A chain is valid if all of the following are true:

1. B0 is a valid genesis block.
2. For every i > 0, Bi.previous\_hash = B(i-1).block\_hash.
3. For every block Bi, recomputed\_hash(Bi) = Bi.block\_hash.
4. For every block Bi, validator signature is valid under the validator identity rules.
5. No duplicate transaction hash violates system policy.

\subsection{Integrity Claim}

If SHA-256 is collision resistant and block verification is performed honestly, then modifying any committed transaction or block field should be detected because the Merkle root or block hash changes. The chain therefore provides tamper evidence, not automatic tamper prevention.

\subsection{Privacy Claim}

If only file hashes and pseudonymous metadata are stored on-chain, the chain does not directly reveal the medical file. However, metadata can still leak sensitive information if patient\_id or doctor\_id are meaningful identifiers. A production system should tokenize or pseudonymize these fields more strongly.

27. DESIGN TRADEOFF ANALYSIS

\subsection{Why Proof of Authority Instead of Proof of Work?}

Proof of Work is designed for open anonymous networks. Healthcare networks are not anonymous; validators are known institutions. Proof of Authority avoids mining cost, reduces latency, and improves accountability. The tradeoff is that governance becomes important. The system must carefully decide who is allowed to validate.

\subsection{Why Off-Chain Storage?}

Off-chain storage protects privacy and reduces blockchain size. The tradeoff is that availability of the original file depends on the off-chain storage system. If the file is deleted, the chain can still prove that a file once existed, but it cannot recover the file.

\subsection{Why C for the Core?}

C gives low-level control, fast execution, and transparent memory layout. It is useful for learning and systems programming. The tradeoff is that C requires careful handling of buffers, string lengths, memory allocation, and platform differences. A production version may prefer Rust or a memory-safe language for reducing memory safety risk.

\subsection{Why Express as a Bridge?}

Express is fast to build and easy to connect with a React frontend. The tradeoff is that the current bridge uses child\_process execution, which is less robust than a persistent core service or native binding. A future version should expose the core through a stable RPC API.

\subsection{Why React for Frontend?}

React is suitable for interactive dashboards and role-based views. It makes the project more demonstrable. The tradeoff is that UI can create a false sense of completeness if backend workflows are not fully implemented. The report must therefore clearly distinguish implemented backend behavior from simulated frontend behavior.

\subsection{Why Direct Binary Ledger Files?}

Direct binary files are easy to append and fast to read in C. The tradeoff is portability. Different platforms may represent structures differently. A canonical serialization format would be safer.

\subsection{Why Demo Signature Mode?}

The current demo signature mode removes OpenSSL dependency and makes the project easier to compile. The tradeoff is security. It must be replaced with real cryptography before any serious deployment.

28. EVALUATION METHODOLOGY FOR FINAL THESIS

The final academic thesis should not only state that the system works. It should evaluate it with repeatable experiments.

\subsection{Functional Tests}

Recommended functional tests:

1. Create genesis block.
2. Add one valid record.
3. Add duplicate record and confirm rejection.
4. Tamper with a block and confirm verification failure.
5. Tamper with off-chain file and confirm hash mismatch.
6. Start two nodes and confirm synchronization.
7. Start four nodes and confirm consensus commit.
8. Query blocks through API and confirm frontend display.

\subsection{Performance Tests}

Recommended metrics:

\begin{itemize}[leftmargin=1.8em]
\item block creation time
\item block verification time
\item file hashing time by file size
\item block append time
\item API response time
\item frontend load time
\item synchronization time for N missing blocks
\item throughput in blocks per second
\item throughput in transactions per second
\end{itemize}

\subsection{Suggested Experimental Matrix}

Run tests with:

\begin{itemize}[leftmargin=1.8em]
\item 1 node, 10 records
\item 1 node, 100 records
\item 1 node, 1000 records
\item 2 nodes, 100 records
\item 4 nodes, 100 records
\item files of 10 KB, 1 MB, and 10 MB
\end{itemize}

Each test should be repeated at least five times. The report should include average, minimum, maximum, and standard deviation.

\subsection{Stress Tests}

Recommended stress tests:

\begin{itemize}[leftmargin=1.8em]
\item repeated duplicate submissions
\item rapid API polling
\item node shutdown during sync
\item node shutdown during proposal
\item corrupted blockchain file
\item missing off-chain record
\item malformed network message
\end{itemize}

\subsection{Security Tests}

Recommended security tests:

\begin{itemize}[leftmargin=1.8em]
\item forged block signature
\item wrong previous hash
\item duplicate transaction hash
\item modified transaction data
\item invalid validator port
\item oversized serialized message
\item path traversal attempt during upload
\end{itemize}

29. RISK REGISTER

Risk 1: Mock signature accepted as real security

Impact: High
Likelihood: Medium
Mitigation: Clearly document demo status and replace with real cryptography.

Risk 2: Frontend upload appears complete but backend does not commit records

Impact: High
Likelihood: High
Mitigation: Implement a real add-record API path and integration test.

Risk 3: Plain TCP exposes node messages

Impact: Medium to High
Likelihood: Medium
Mitigation: Add TLS or mutual TLS.

Risk 4: Binary ledger portability problems

Impact: Medium
Likelihood: Medium
Mitigation: Use canonical serialization.

Risk 5: Lost off-chain files

Impact: High
Likelihood: Medium
Mitigation: Use redundant encrypted storage, backups, and file availability checks.

Risk 6: Weak patient privacy through metadata

Impact: High
Likelihood: Medium
Mitigation: Pseudonymize patient and doctor identifiers; minimize metadata.

Risk 7: Consensus failure during network partition

Impact: High
Likelihood: Medium
Mitigation: Formalize consensus, quorums, timeouts, and recovery.

Risk 8: Child process bridge injection or operational fragility

Impact: Medium
Likelihood: Medium
Mitigation: Avoid shell command composition; expose core through stable service API.

Risk 9: No automated regression tests

Impact: Medium
Likelihood: High
Mitigation: Add unit, integration, and multi-node tests.

Risk 10: Repository contains generated data and binaries

Impact: Low to Medium
Likelihood: High
Mitigation: Clean .gitignore and separate demo artifacts from source code.
\textbackslash\{\}n\textbackslash\{\}n
\subsubsection{References}

~\cite{1} Uncovering Security Vulnerabilities in Electronic Medical Record Systems: A Comprehensive Review of Threats and Recommendations for Enhancement, accessed on June 18, 2026, https://eprints.uad.ac.id/61866/1/6-Uncovering\%20Security\%20Vulnerabilities\%20in\%20Electronic\%20Medical\%20Record\%20Systems\%20A\%20Comprehensive\%20Review\%20of\%20Threats\%20and\%20Recommendations\%20for\%20Enhancement.pdf
A blockchain-enabled sharing platform for personal health records - PMC, accessed on June 18, 2026, https://pmc.ncbi.nlm.nih.gov/articles/PMC10366433/
A multi-layered cryptographic trust reinforcement model against AI-driven threat propagation and zero-day cloud vulnerabilities in healthcare data ecosystems - PMC, accessed on June 18, 2026, https://pmc.ncbi.nlm.nih.gov/articles/PMC12920886/
updated\_thesis\_report.txt
A blockchain-based verifiable CP-ABE scheme for medical data privacy protection - PMC, accessed on June 18, 2026, https://pmc.ncbi.nlm.nih.gov/articles/PMC12301491/
A Blockchain-Based Electronic Health Record (EHR) System for ..., accessed on June 18, 2026, https://www.mdpi.com/2073-431X/13/6/132
The legal considerations of AI-blockchain for securing health data - NCBI - NIH, accessed on June 18, 2026, https://www.ncbi.nlm.nih.gov/books/NBK613198/
Securing healthcare data with blockchain in 2025 - Paubox, accessed on June 18, 2026, https://www.paubox.com/blog/securing-healthcare-data-with-blockchain-in-2025
Secure, scalable, and interoperable healthcare data exchange using layer-2 ZK-rollups, smart contracts, and IPFS - PMC, accessed on June 18, 2026, https://pmc.ncbi.nlm.nih.gov/articles/PMC12901292/
Blockchain-Based Secure Data Sharing In Healthcare: Concepts, Applications, And Research Directions - XLESCIENCE, accessed on June 18, 2026, https://xlescience.org/index.php/IJASIS/article/download/880/324
Blockchain Bridges Critical National ... - City Research Online, accessed on June 18, 2026, https://openaccess.city.ac.uk/id/eprint/32135/
Securing electronic health records against insider-threats: A supervised machine learning approach - LJMU Research Online, accessed on June 18, 2026, https://researchonline.ljmu.ac.uk/id/eprint/18592/1/Securing\%20electronic\%20health\%20records\%20against\%20insider-threats\%20A\%20supervised\%20machine\%20learning\%20approach.pdf
Advanced Privacy-Preserving Decentralized Federated Learning for Insider Threat Detection in Collaborative Healthcare Institutions - ProQuest, accessed on June 18, 2026, https://search.proquest.com/openview/e8a68ee960829e5def9299058a2c6a3c/1?pq-origsite=gscholar\&cbl=18750\&diss=y
A Comprehensive Survey of Cybersecurity Threats and Data Privacy Issues in Healthcare Systems - MDPI, accessed on June 18, 2026, https://www.mdpi.com/2076-3417/16/3/1511
A multi-layered cryptographic trust reinforcement model against AI-driven threat propagation and zero-day cloud vulnerabilities in healthcare data ecosystems - PubMed, accessed on June 18, 2026, https://pubmed.ncbi.nlm.nih.gov/41634110/
Shitanshu jain - Google Scholar, accessed on June 18, 2026, https://scholar.google.com/citations?user=9t2euX8AAAAJ\&hl=en
A lightweight authentication and authorization method in IoT-based medical care, accessed on June 18, 2026, https://www.semanticscholar.org/paper/A-lightweight-authentication-and-authorization-in-Khajehzadeh-Barati/1ce9d4bb32c33d709922a2813e061ea8ac263574
CMES | Free Full-Text | Security and Privacy Challenges, Solutions, and Performance Evaluation in AIoT-Enabled Smart Societies - Tech Science Press, accessed on June 18, 2026, https://www.techscience.com/CMES/v146n3/66792/html
A Case Study for Blockchain in Healthcare : “ MedRec ” prototype for electronic health records and medical research data - Semantic Scholar, accessed on June 18, 2026, https://www.semanticscholar.org/paper/A-Case-Study-for-Blockchain-in-Healthcare-\%3A-\%E2\%80\%9C-\%E2\%80\%9D-for-Ekblaw-Azaria/3ed0db58a7aec7bafc2aa14ca550031b9f7021d5
MedRec: Using Blockchain for Medical Data Access and Permission Management, accessed on June 18, 2026, https://people.cs.pitt.edu/~babay/courses/cs3551/papers/MedRec.pdf
MedRec: Using Blockchain for Medical Data Access and Permission Management, accessed on June 18, 2026, https://www.researchgate.net/publication/308570159\_MedRec\_Using\_Blockchain\_for\_Medical\_Data\_Access\_and\_Permission\_Management
Blockchain Technology for Electronic Health Records - PMC - NIH, accessed on June 18, 2026, https://pmc.ncbi.nlm.nih.gov/articles/PMC9739765/
Ancile: Privacy-Preserving Framework for Access Control and Interoperability of Electronic Health Records Using Blockchain Technology - ScholarWorks, accessed on June 18, 2026, https://scholarworks.boisestate.edu/cs\_facpubs/140/
Ancile: Privacy-preserving Framework for Access Control and Interoperability of Electronic Health Records Using Blockchain Technology - ResearchGate, accessed on June 18, 2026, https://www.researchgate.net/publication/323247610\_Ancile\_Privacy-preserving\_Framework\_for\_Access\_Control\_and\_Interoperability\_of\_Electronic\_Health\_Records\_Using\_Blockchain\_Technology
Blockchain Personal Health Records: Systematic Review - PMC - NIH, accessed on June 18, 2026, https://pmc.ncbi.nlm.nih.gov/articles/PMC8080150/
FHIRChain: Applying Blockchain to Securely and Scalably Share Clinical Data - PMC, accessed on June 18, 2026, https://pmc.ncbi.nlm.nih.gov/articles/PMC6082774/
FHIRChain: Applying Blockchain to Securely and Scalably Share Clinical Data, accessed on June 18, 2026, https://www.dre.vanderbilt.edu/~schmidt/PDF/FHIRChain-jnca.pdf
A Quantitative Landscape Analysis of Blockchain for EHR Security and Interoperability, accessed on June 18, 2026, https://www.isjtrend.com/article\_227487.html
Biomedical blockchain with practical implementations and quantitative evaluations: a systematic review - Oxford Academic, accessed on June 18, 2026, https://academic.oup.com/jamia/article/31/6/1423/7668012
Secure electronic health record access control via blockchain, dual-attribute encryption, and large language model-based attribute extraction - PMC, accessed on June 18, 2026, https://pmc.ncbi.nlm.nih.gov/articles/PMC12979818/
Privacy preservation in blockchain-based healthcare data sharing: A systematic review, accessed on June 18, 2026, https://pmc.ncbi.nlm.nih.gov/articles/PMC12534302/
Dual blockchain-based data sharing mechanism with privacy protection for medical internet of things - PMC, accessed on June 18, 2026, https://pmc.ncbi.nlm.nih.gov/articles/PMC10758875/
MeDShare: Trust-Less Medical Data Sharing Among Cloud Service Providers via Blockchain | Semantic Scholar, accessed on June 18, 2026, https://www.semanticscholar.org/paper/MeDShare\%3A-Trust-Less-Medical-Data-Sharing-Among-via-Xia-Sifah/49af9119c09b97af977595b011afd8a3f588412d
(PDF) MeDShare: Trust-less Medical Data Sharing Among Cloud Service Providers Via Blockchain - ResearchGate, accessed on June 18, 2026, https://www.researchgate.net/publication/318666875\_MeDShare\_Trust-less\_Medical\_Data\_Sharing\_Among\_Cloud\_Service\_Providers\_Via\_Blockchain
Analyzing the performance of a blockchain-based personal health record implementation, accessed on June 18, 2026, https://pubmed.ncbi.nlm.nih.gov/30844481/
OmniPHR : A Blockchain based Interoperable Architecture for Personal Health Records | Request PDF - ResearchGate, accessed on June 18, 2026, https://www.researchgate.net/publication/337910407\_OmniPHR\_A\_Blockchain\_based\_Interoperable\_Architecture\_for\_Personal\_Health\_Records
Secure and Trustable Electronic Medical Records Sharing using Blockchain - PMC - NIH, accessed on June 18, 2026, https://pmc.ncbi.nlm.nih.gov/articles/PMC5977675/
Secure and Trustable Electronic Medical Records Sharing using Blockchain - Jamie Munro, accessed on June 18, 2026, https://portfolio.jbm.fyi/secure-and-trustable-electronic-medical-records-sharing-using-blockchain/
Addressing the Challenges of Electronic Health Records Using Blockchain and IPFS - PMC, accessed on June 18, 2026, https://pmc.ncbi.nlm.nih.gov/articles/PMC9183171/
Addressing the Challenges of Electronic Health Records Using Blockchain and IPFS, accessed on June 18, 2026, https://pubmed.ncbi.nlm.nih.gov/35684652/
(PDF) A Secure and Interoperable Architecture for Blockchain/IPFS Assisted Electronic Health Record Access Control and Sharing - ResearchGate, accessed on June 18, 2026, https://www.researchgate.net/publication/373005704\_A\_Secure\_and\_Interoperable\_Architecture\_for\_BlockchainIPFS\_Assisted\_Electronic\_Health\_Record\_Access\_Control\_and\_Sharing
Electronic Health Record Systems Based on Blockchain: A Comprehensive Survey - MDPI, accessed on June 18, 2026, https://www.mdpi.com/2076-3417/16/8/3768
A Secure and Interoperable Architecture for Electronic Health Record Access Control and Sharing - arXiv, accessed on June 18, 2026, https://arxiv.org/pdf/2602.20830
[Article] HealthRec-Chain: Patient-centric blockchain enabled IPFS for privacy preserving scalable health data - Reddit, accessed on June 18, 2026, https://www.reddit.com/r/ipfs/comments/1tvr8pl/article\_healthrecchain\_patientcentric\_blockchain/
[2503.00742] Revolutionizing Healthcare Record Management: Secure Documentation Storage and Access through Advanced Blockchain Solutions - arXiv, accessed on June 18, 2026, https://arxiv.org/abs/2503.00742
Toward blockchain based electronic health record management with fine grained attribute based encryption and decentralized storage mechanisms - PMC, accessed on June 18, 2026, https://pmc.ncbi.nlm.nih.gov/articles/PMC12494854/
A Secure and Interoperable Architecture for Electronic Health Record Access Control and Sharing - arXiv, accessed on June 18, 2026, https://arxiv.org/html/2602.20830v1
MedBlock: Secure Medical Data Sharing | PDF | Electronic Health Record - Scribd, accessed on June 18, 2026, https://www.scribd.com/document/465212996/p74-MedBlock-Efficient-and-Secure-Medical-Data-Sharing-Via-Blockch
Security research of blockchain technology in electronic medical records - PMC, accessed on June 18, 2026, https://pmc.ncbi.nlm.nih.gov/articles/PMC9439767/
View of Blockchain Technology for Enhancing Clinical Data Management in Healthcare: A Systematic Literature Review, accessed on June 18, 2026, https://blockchainhealthcaretoday.com/index.php/journal/article/view/471/1032
Analyzing the Prospects of Blockchain in Healthcare Industry - PMC, accessed on June 18, 2026, https://pmc.ncbi.nlm.nih.gov/articles/PMC9733997/
The role of blockchain as a security mechanism for electronic health records - Diva-Portal.org, accessed on June 18, 2026, https://www.diva-portal.org/smash/get/diva2:1879894/FULLTEXT01.pdf
Enhancing Network Lifetime in e-Healthcare Systems using a Proof-of-Authority Permissioned Blockchain Approach - IEEE Xplore, accessed on June 18, 2026, https://ieeexplore.ieee.org/document/11200952/
View of Design and Development of a Blockchain-Based System to Protect Medical Records, accessed on June 18, 2026, https://blockchainhealthcaretoday.com/index.php/journal/article/view/381/1040
A blockchain framework using proof of authority and smart contracts for ethical and secure healthcare asset management - PMC, accessed on June 18, 2026, https://pmc.ncbi.nlm.nih.gov/articles/PMC12500698/
Patient Identity Protection and Duplicate Record Prevention in Electronic Health Record (EHR) Systems - IEEE Xplore, accessed on June 18, 2026, https://ieeexplore.ieee.org/document/11431915/
A Survey of AI-Driven Techniques for Enhancing Clinical Decision Support in Modern Healthcare Systems | International Journal of Emerging Research in Engineering and Technology, accessed on June 18, 2026, https://ijeret.org/index.php/ijeret/article/view/619
Eliminating Duplicate Medical Records: How Modern Solutions are Revolutionizing Healthcare Data Management - ResearchGate, accessed on June 18, 2026, https://www.researchgate.net/publication/391159020\_Eliminating\_Duplicate\_Medical\_Records\_How\_Modern\_Solutions\_are\_Revolutionizing\_Healthcare\_Data\_Management
Chandra Sekhara Reddy Adapa's research works - ResearchGate, accessed on June 18, 2026, https://www.researchgate.net/scientific-contributions/Chandra-Sekhara-Reddy-Adapa-2342938594
Designing and testing a blockchain application for patient identity management in healthcare - PMC, accessed on June 18, 2026, https://pmc.ncbi.nlm.nih.gov/articles/PMC7928860/\textbackslash\{\}n
\end{document}
